% Workshop version for NeurIPS 2026 "Interpretability as a Science" (InterpScience).
% 9-page long-paper track; references and appendices do not count toward the limit.
% Condensed from mechanistic_validity_v23.tex; audits ride along as Appendix A.
\pdfoutput=1

\documentclass{article}

% Submissions accept ICLR or NeurIPS format; camera-ready must be NeurIPS.
% Swap in neurips_2026.sty when the 2026 kit is released.
\usepackage[preprint]{neurips_2025}
% the style's [preprint] footline claims "Under review."; nothing is submitted yet
\makeatletter\renewcommand{\@noticestring}{Preprint.}\makeatother

\usepackage{amsmath,amssymb}
\usepackage{booktabs}
\usepackage{graphicx}
\graphicspath{{figures/}}
\usepackage{xcolor}
\usepackage{hyperref}
\usepackage{cleveref}
\usepackage{enumitem}
\usepackage{placeins}
\usepackage{float}
\usepackage{multirow}
\usepackage{array}
\usepackage{longtable}

\usepackage{caption}
\captionsetup[table]{skip=8pt}

% let text share a page with the large framework tables
\renewcommand{\topfraction}{0.95}
\renewcommand{\bottomfraction}{0.6}
\renewcommand{\textfraction}{0.05}
\renewcommand{\floatpagefraction}{0.92}
\setcounter{topnumber}{3}
\setcounter{totalnumber}{4}

% tight tables
\newcommand{\tighttable}{\setlength{\tabcolsep}{4pt}\renewcommand{\arraystretch}{1.05}}

\title{Mechanistic Validity: A Validity Theory and Evidence Standard for Mechanistic Claims}

\author{%
  Elliot Tower \\
  Independent Researcher \\
  \texttt{elliot@elliottower.ai}
}

% each audited claim opens at the top of its own page, as in the supplement
\newcommand{\claimheading}[1]{\clearpage\subsection{#1}}

\begin{document}

\maketitle

\begin{abstract}
A mechanistic claim is only as strong as its weakest validity dimension. Five validity types---construct, measurement, internal, external, and interpretive---form a dependency chain: construct validity is a ceiling on all others, so no amount of causal evidence rescues a claim whose target concept is incoherent. We operationalize this theory as Mechanistic Validity, 36 falsifiable criteria adapted from philosophy of science, measurement theory, and validation methodology in neuroscience, pharmacology, and genetics, and we audit sixteen mechanistic claims drawn from fifteen papers. The framework discriminates, and it separates claims that share a paper: induction heads reach Triangulated for in-context token copying and Disconfirmed for general in-context learning, on evidence the original authors named as their own refutation condition. The IOI circuit reaches Causally Suggestive under its most lenient reading, capped by method disagreement: its headline faithfulness number spans below 0\% to over 100\% depending on the ablation protocol. We propose the audit format used here as a standard reporting unit for mechanistic claims: a scorecard over the criteria, the readings the claim admits, and the tier each reading reaches.
\end{abstract}

% ============================================================
\section{Introduction}
\label{sec:intro}

Mechanistic interpretability has produced a remarkable catalog of internal mechanisms:
induction heads that implement in-context copying \citep{olsson2022context}, a multi-role
circuit for indirect object identification \citep{wang2023interpretability}, a greater-than
algorithm \citep{hanna2023greater}, and sparse autoencoder features proposed to correspond
to human-interpretable concepts \citep{cunningham2024sparse, bricken2023monosemanticity}, alongside growing evaluation infrastructure \citep{miller2024transformer, mueller2025mib, karvonen2025saebench, chan2022causal}.

But how well-supported are these claims? Across the field's most prominent results, a pattern recurs: strong evidence along one dimension, untested assumptions along others. Faithfulness scores that are method-conditional. Metrics that saturate on random models. Specificity controls that were never run. Labels that imply more than the evidence supports. Each is a gap in the \emph{kind} of evidence collected, which no amount of execution quality would close, and the pattern repeats in every field that makes mechanistic claims (\Cref{sec:motivation}). Shared metrics are not yet shared validity criteria.

Mechanistic Validity is a structured framework that takes the validity vocabulary the field already uses informally and makes it operational, falsifiable, and auditable. It evaluates claims through a seven-layer pipeline (\Cref{sec:framework}) running from a declared \emph{description mode} through evidence families, metrics, criteria, and five \emph{validity types}---construct, measurement, internal, external, and interpretive, in dependency order---to a \emph{verdict}: the structured record of which criteria are met, together with the highest tier the evidence reaches, from Proposed through Causally Suggestive, Mechanistically Supported, and Triangulated, to Validated.

The pipeline runs in both directions: evidence builds upward from metrics through criteria to a verdict, and an audit reverses the process, identifying which criteria are unmet and whether the description mode is supported. The same criteria apply to circuits, SAE features, probing classifiers, steering vectors, transcoders, and crosscoders---the questions ``is the construct well-defined?'' and ``is the evidence causal?'' do not depend on how the evidence was produced. The unit of analysis is a single causal mechanistic claim, not a model, a method, or a benchmark score. The full 36-criterion audit of each claim appears in \Cref{appendix:audits}; an extended version carries the cross-field failure catalog, the theoretical foundations, and the 59-metric inventory in full \textcolor{red}{[CITE: extended version]}.

% ============================================================
\section{Validity Failures Across Sciences}
\label{sec:motivation}

Every field that makes mechanistic claims has encountered the same structural problem: strong individual results that do not add up to valid conclusions. \citet{bennett2010neural} scanned a dead Atlantic salmon while presenting it with photographs of humans in social situations; at standard uncorrected thresholds, 16 significant voxels appeared in the salmon's brain cavity---a missing baseline separation (M2). For three decades stenting was standard treatment for stable angina; the first sham-controlled trial found no significant difference in symptoms, so the causal attribution had never been tested against placebo \citep{alhabib2018orbita}---a missing specificity control (I4). And ``emergent abilities'' of large language models are artifacts of nonlinear evaluation metrics: switch exact-match accuracy for a linear metric and the phase transition disappears \citep{schaeffer2023emergent}---a calibration failure (M4).

Mechanistic interpretability already exhibits the same structures. \citet{karvonen2025saebench} run a random-initialization control on five of their eight SAE evaluation metrics, and one scores \emph{higher} on a randomly initialized model than on a trained one; unconstrained nonlinear alignment maps achieve 100\% interchange-intervention accuracy on randomly initialized models \citep{sutter2025nonlinear}; and automated interpretability scores on sparse autoencoders trained over randomly initialized transformers are similar to those from trained models \citep{heap2025sparse}. One control---baseline separation---catches all three, and it tests whether a reported quantity is about the model at all. In each of these failures, across fields, the individual experimental result was strong, but the inference chain from measurement to conclusion had an untested link: claims migrated upward in inferential scope faster than supporting validity criteria could be established. \Cref{tab:cross-science} in \Cref{appendix:cross-science} maps each failure to the criterion that detects it.

% ============================================================
\section{Related Work}
\label{sec:related}

Validity frameworks in wide use fix the proposition under assessment before assessing it, because their fields supply one. \citet{cronbach1955construct} and \citet{messick1995validity} developed construct validity for measurement instruments; \citeauthor{campbell1963experimental}'s taxonomy, extended by \citet{shadish2002experimental}, evaluates causal inference from experiments; \citet{guyatt2008grade} rate the certainty of a body of evidence bearing on one clinical question. None of these objects is a claim about how a system works. A mechanistic claim has parts, an organization, a level of description, and a scope, and each can be asserted at more or less strength independently of the others. Those four dimensions are why the framework needs machinery the others do not: description modes, an interpretive validity block, and a procedure for identifying which proposition is being assessed before assessing it. Recent work applies construct validity to ML evaluation more broadly \citep{bean2025measuring, freiesleben2025benchmarking}; \citeauthor{freiesleben2025benchmarking} hold that their scheme implies no strict hierarchy between validity types, where this framework orders the five types by dependency.

The observation that claims carry variable strength is old \citep{toulmin1958uses}, and four recent proposals arrange one claim by inferential strength \citep{leavitt2020falsifiable, gupta2026anthropomorphic, salaudeen2025measurement, jacovi2020faithfully}. The nearest precedent is the surrogate-outcome literature in evidence-based medicine, which scores one relationship against three explicitly separate levels of evidence, where satisfying a weaker level establishes nothing about a stronger one \citep{euhta2024outcomes}. That is the structure this paper generalizes. What is distinctive here is that readings are derived from what a given paper asserted rather than fixed by domain, vary along two stated axes---mechanistic specificity and scope---and are all scored against one evidence record, so the strongest supported reading can be reported without erasing the stronger formulations the paper or its citers use elsewhere. The stronger formulations matter because of \emph{citation transmutation}: a hypothesis becomes an asserted fact through citation alone, producing the gap between what a paper claimed and what the field takes it to have shown \citep{greenberg2009citation}.

Existing interpretability infrastructure evaluates complementary parts of the pipeline. MIB \citep{mueller2025mib} scores localization methods against each other on fixed tasks; Tracr \citep{lindner2023interp} compiles programs into weights, supplying ground truth; causal scrubbing \citep{chan2022causal} and the hypothesis tests of \citet{shi2024hypothesis} take a supplied hypothesis and ask whether it survives. Mechanistic Validity takes the output of these pipelines as its input: given a claim already published, it asks which readings of that claim the whole evidence record supports. SAEBench \citep{karvonen2025saebench} audits SAE metrics, and this framework treats that audit as evidence about its own measurement criteria. \citet{sharkey2025open} count the validation of interpretability findings among the field's open problems; \citet{orgad2026actionable} propose judging interpretability by actionability---a different standard, since a finding can be actionable and unwarranted. A companion work \citep{tower2026mechviews} sets out nine views of what a mechanism is; this framework takes a claim's view as an input, answering which rival within that view the evidence favors and how strong a reading the winner will carry.

% ============================================================
\section{The Framework}
\label{sec:framework}

The framework evaluates a mechanistic claim through seven layers, in pipeline order, each answering one question and taking the previous layer's answer as input:
\begin{enumerate}[nosep,leftmargin=*]
\item \textbf{Description mode.} What kind of explanation is offered?
\item \textbf{Evidence families.} What sources of signal support it?
\item \textbf{Metrics.} What was concretely measured?
\item \textbf{Criteria.} Does the evidence meet the stated conditions?
\item \textbf{Validity types.} Which dimensions of validity does it address?
\item \textbf{Synthesis.} How is evidence aggregated across methods?
\item \textbf{Verdict.} What has the claim established?
\end{enumerate}
The framework defines the structure of evaluation; a companion methodology defines the concrete procedures for generating evidence---59 metrics, including the 15 calibrations that decide measurement verdicts. Layer 6 has no procedure to describe: where several discovery methods return overlapping but non-identical component sets, nothing in current practice adjudicates between them, and every audit in \Cref{sec:case-studies} records that gap rather than closing it.

\subsection{Description modes}
\label{sec:description-modes}

A description mode is a commitment about what kind of explanation is being offered. The declaration comes before the analysis and determines what counts as evidence, what counts as a gap, and what the finding means if it succeeds. The seven modes extend Marr's three-level hierarchy \citep{marr1982vision}---computational (what function the system computes, and why), algorithmic (what procedure it executes), representational (what information is encoded, in what geometry)---by splitting the implementational level into four sub-types: functional (what transformation each component performs), connectomic (how components are wired), statistical (what activations look like), and topographic (which components are involved). \Cref{tab:modes} in \Cref{appendix:modes} gives an example claim at each mode. The modes form a partial order by evidential commitment, and a higher mode requires the evidence of the modes below it plus bridging evidence; the verdict is issued at the strongest mode whose required evidence is complete, with partial higher-mode evidence reported as suggestive rather than counted toward an upgrade.

Mechanistic interpretability routinely conflates levels, and Anthropic's global-workspace result illustrates the conflation: the Jacobian lens recovers a residual-stream subspace (implementational-statistical), the paper says that subspace encodes what the model can report (representational), and it narrates the whole as a ``global workspace'' the model reasons in (computational)---three modes deep, on evidence that reaches the representational level at best \citep{gurnee2026verbalizable}. Declaring the mode first forces the question the label skips: what bridging evidence carries a subspace claim up to a computational one?

\subsection{Evidence families}
\label{sec:evidence-families}

An evidence family classifies a metric's output by the source of signal it draws on. Evidence for a mechanism claim falls into four families: \textbf{structure}, the system's persistent organization, observable without running it; \textbf{state}, what the system does internally while operating; \textbf{behavior}, its input-output relation; and \textbf{history}, how it came to be as it is. In mechanistic interpretability these are weights, activations, task performance, and training. Each family crosses with a second axis---observational vs.\ interventional---producing eight cells, and most published evidence lives in a single cell: interventional activations, meaning patching and ablation. Triangulation means covering multiple cells, and the observational/interventional axis matters as much as the family axis: two interventional-activation metrics share confounds that an observational-weight metric does not.

\subsection{Criteria}
\label{sec:criteria}

Criteria are specific, falsifiable conditions that must be met for a validity type to be satisfied (\Cref{tab:criteria}). Each criterion has a concrete pass condition---necessity (I1) requires that ablating the component degrades performance across at least two independent methods---and each receives one of six statuses: Confirmed, Partially confirmed, Inconclusive, Untested, Disconfirmed, Not applicable. Confirmed and Partially confirmed establish a criterion. Inconclusive, Untested and Disconfirmed do not, and tier promotion sees no difference among them: a criterion that was tested and failed and one that nobody tested both leave the criterion unestablished, and ranking them would make not running a test preferable to running one and reporting an adverse result. Not applicable is argued from what the claim is, never from what an experiment would cost, and is excluded from the denominator rather than counted as passing. Statuses are categorical evidential states, not posterior beliefs \citep{mayo2018statistical}; recording insufficiency as a graded conclusion follows the US Preventive Services Task Force's \emph{I statement} \citep{uspstf2021}, whose single category we divide in two because the remedies differ: an Untested criterion needs an experiment, an Inconclusive one needs adjudication between results that already exist.

\begingroup
\hyphenpenalty=10000
\exhyphenpenalty=10000
\begin{table}[tbp]
\centering
\footnotesize
\tighttable
\caption{Thirty-six criteria across five validity types. Each criterion has a concrete test, and each receives one of six statuses, producing a structured validity profile.}
\label{tab:criteria}
\begin{tabular}{@{}llp{7.6cm}@{}}
\toprule
\textbf{ID} & \textbf{Criterion} & \textbf{Description} \\
\midrule
C1 & Falsifiability & Can the claim be refuted? \\
C2 & Structural plausibility & Is the mechanism physically possible in the architecture? \\
C3 & Convergent validity & Do multiple independent methods agree? \\
C4 & Discriminant validity & Does the measure distinguish this from neighboring constructs? \\
C5 & Nomological validity & Does the claim fit into a broader theory? \\
C6 & Complementation validity & Are the construct's labeled subdivisions functionally distinct? \\
\midrule
M1 & Reliability & Do repeated measurements give the same answer? \\
M2 & Baseline separation & Is the score distinguishable from random/untrained baselines? \\
M3 & Stability & Is the classification robust to perturbation? \\
M4 & Calibration & Are the numbers meaningful? \\
M5 & Sensitivity & Can the instrument detect known-true effects? \\
M6 & Invariance & Does the metric behave consistently across conditions? \\
M7 & Selection correction & When $k$ findings are selected from $N$ candidates, is $N$ reported and multiplicity controlled? \\
\midrule
I1 & Necessity & Is the circuit required for the behavior? \\
I2 & Sufficiency & Is the circuit enough to produce the behavior? \\
I3 & Minimality & Does every component earn its place? \\
I4 & Specificity & Does intervening on the circuit affect this task more than matched control tasks? \\
I5 & Rival mechanism exclusion & Is this \emph{the} mechanism, or \emph{a} mechanism? \\
I6 & Double dissociation & Do two interventions cross, each breaking what the other spares? \\
I7 & Confound control & Are alternative explanations ruled out? \\
I8 & Confounding sensitivity & How strong must an unmeasured confounder be to explain the result? \\
I9 & Epistatic interaction & Do circuit components interact non-additively, and does the direction distinguish shared pathways from mutual compensation? \\
I10 & Rescue reversibility & Does restoring a corrupted component recover behavior? \\
I11 & Onset coupling & Does the mechanism appear when the capability appears? \\
I12 & Offset coupling & Does the mechanism go when the capability is removed? \\
\midrule
E1 & Intervention reach & Has the result been reproduced under at least two intervention families, and do they agree? \\
E2 & Prompt generalization & Does it work on diverse prompts? \\
E3 & Cross-task generalization & Does the mechanism transfer to related tasks? \\
E4 & Cross-model generalization & Does the mechanism appear in other models? \\
E5 & Graded response & Does partial ablation produce partial effects? \\
E6 & Novel prediction & Does the mechanism predict new, untested behaviors? \\
\midrule
V1 & Level declaration & At what description mode is the claim made? \\
V2 & Level-evidence match & Does the evidence support claims at that level? \\
V3 & Alternative level & Could the evidence be explained at a different level? \\
V4 & Unlicensed labeling & Does a name import a property that was not measured? \\
V5 & Scope declaration & What does the claim explicitly not cover? \\
\bottomrule
\end{tabular}
\end{table}
\endgroup

\subsection{Validity types}
\label{sec:validity-types}

The five validity types form a dependency chain:
\begin{equation}
\text{Construct} \rightarrow \text{Measurement} \rightarrow \text{Internal} \rightarrow \text{External} \rightarrow \text{Interpretive}
\label{eq:dependency}
\end{equation}
The ordering reflects logical precedence: a construct that is not well-defined cannot be reliably measured, an unreliable measurement cannot support a causal claim, a causal claim in one setting cannot be generalized, and an interpretation cannot be evaluated until the underlying evidence is characterized. The vocabulary is \citeauthor{shadish2002experimental}'s; the ordering is not. We place construct first, following \citet{messick1995validity}, for whom construct validity ``undergirds all score-based interpretations.'' Reading the remaining four as a strict chain is our own commitment, and the audits are what test it: if a claim could establish external validity while failing measurement validity, the chain would be wrong. The chain is a ceiling, not a sequence of gates.

Construct validity asks whether the target concept is well-defined---falsifiable, structurally plausible, and distinguishable from neighboring constructs: a ``deception feature'' that cannot be distinguished from an ``uncertainty feature'' fails regardless of how carefully it is measured. Measurement validity asks whether the instruments are trustworthy: reliable, separable from random baselines, and invariant across conditions \citep{borsboom2004concept}. Internal validity asks whether the evidence establishes that the component is causally involved---necessity, sufficiency, and specificity, with double dissociation rather than single dissociation as the test that separates an implementation from a bottleneck \citep{shallice1988neuropsychology, craver2007explaining, pearl2009causality}. External validity asks whether the mechanism generalizes across prompts, ablation methods, tasks, and models \citep{pearl2011transportability}. Interpretive validity asks whether the natural-language description is justified: declaring the Marr level at which the claim operates, matching evidence to that level, and checking that the label does not import a property never measured \citep{marr1982vision, messick1995validity}.

Two scoring rules follow from the chain. Reliability (M1) bounds internal validity in a strict sense, since attenuation caps the correlation two measures can exhibit at the geometric mean of their reliabilities; stability (M3) and invariance (M6) instead record that a reported quantity moves with analysis choices, which impugns the published number without ceilinging the effect it estimates, so M1 caps and M3/M6 are reported as qualifications. And a double dissociation (I6) is a single crossed design whose halves are already scored elsewhere (I1 and M2), so it is met or unmet with no partial credit.

\subsection{Verdict tiers}
\label{sec:verdict-tiers}

Claims are assigned to one of five verdict tiers representing qualitative transitions in evidential status (\Cref{tab:verdicts}), implementing the epistemic reading of \citeauthor{craver2006when}'s how-possibly to how-actually dimension. We report the tier together with the criteria that cap it, as \citet{guyatt2008grade} report a certainty rating alongside the domains that lowered it. The verdict tells researchers exactly what additional evidence would advance the claim: a Proposed claim needs causal intervention (I1) to reach Causally Suggestive; a Causally Suggestive claim needs sufficiency (I2) and method consistency (E1) to reach Mechanistically Supported.

\begin{table}[tbp]
\centering
\footnotesize
\tighttable
\caption{Verdict tiers with minimum evidence requirements. Each tier subsumes all requirements of lower tiers. Underdetermined, Insufficient and Disconfirmed replace the tier rather than accompanying it.}
\label{tab:verdicts}
\begin{tabular}{@{}>{\raggedright\arraybackslash}p{2.1cm}>{\raggedright\arraybackslash}p{3.2cm}>{\raggedright\arraybackslash}p{6.9cm}@{}}
\toprule
\textbf{Tier} & \textbf{Meaning} & \textbf{Minimum evidence requirements} \\
\midrule
Proposed & Structural or representational evidence only & Construct defined under a declared description mode (C1--C2); at least one admissible measurement conducted \\
\midrule
Causally Suggestive & Necessity shown, sufficiency not established & Necessity via causal intervention (I1 confirmed); at least one measurement passes baseline separation (M2) \\
\midrule
Mechanistically Supported & Necessity + sufficiency with consistent methods & Sufficiency established (I2); intervention reach across $\geq 2$ ablation methods (E1); specificity test conducted (I4 at least partially confirmed) \\
\midrule
Triangulated & Multiple converging lines of independent evidence & Convergent evidence from $\geq 3$ evidence families with independent failure modes (C3); cross-distribution replication (E4 where the claim asserts reach beyond the systems tested, E2 where it does not); double dissociation attempted (I6) \\
\midrule
Validated & All five validity types addressed & Measurement calibration (M1--M6) explicitly addressed; interpretive validity (V1--V5) audited; external validity across models (E4) and prompts (E2) \\
\midrule
Underdetermined & Evidence consistent with multiple mechanisms & Cannot resolve between rival specifications \\
\midrule
Insufficient & Cannot be assessed & The construct is not defined well enough, or no admissible measurement exists, to score the claim \\
\midrule
Disconfirmed & Fails decisively on a key criterion & Negative result on a required criterion \\
\bottomrule
\end{tabular}
\end{table}

Some criteria cannot be tested for some claims, and the tier system treats that as information rather than as an exemption: specificity (I4) is unanswerable for a model trained on a single task, onset coupling (I11) for a released model without training checkpoints, and in each case the claim caps at the tier below---a statement about the setting rather than a deficiency in the work. Three diagnostic labels sit outside the hierarchy, covering the three ways a claim can fail to occupy a tier: too many explanations fit it (Underdetermined), too little evidence exists to assess it (Insufficient), or it has been tested and failed (Disconfirmed).

The top tiers require more than accumulated confirmations. Following the error-statistical tradition \citep{mayo2018statistical}, the criteria that lift a claim to Triangulated or Validated must be established by \emph{severe} tests---tests the claim could have failed and did not. Convergent evidence across families with different failure modes (C3) is, by construction, evidence a claim was not optimized to produce, so a claim resting on one family stops at Mechanistically Supported however often it is re-tested. Preregistration is recorded and never scored: scoring it would make identical evidence count differently according to whether a disclosure was made, and this framework grades evidence. Finally, a verdict states how well a claim is evidenced, not how much of the mechanism it specifies; the two vary independently \citep{craver2006when}, mechanistic specificity is bought with evidential risk, and we report both.

% ============================================================
\section{Audits of Sixteen Claims}
\label{sec:case-studies}

We apply the framework to sixteen published mechanistic claims, drawn from fifteen papers (\Cref{tab:verdict-summary}). We audit \emph{claims}, not papers: a verdict characterizes the current evidential status of a claim, aggregating all published evidence to date. Because verdicts are graded against an idealized standard, a low tier marks distance from that ideal, not a flawed study. Every claim was re-derived against all 36 criteria from the origin paper read in full, with the post-origin literature checked against primary sources. Where a criterion is Untested we distinguish \emph{unattempted} from \emph{not applicable}, because only the first is a gap the authors could have closed. Where the field's reading is stronger than the origin paper's claim, the audit says which is being scored. The sixteen claims spread across five labels, with each difference traceable to a named criterion; the full scorecards, readings tables, and tier ladders appear in \Cref{appendix:audits}.

\begingroup
\hyphenpenalty=10000
\begin{table}[tbp]
\centering
\scriptsize
\tighttable
\caption{Verdicts for the sixteen audited claims, at each claim's primary reading. The last column names the criteria blocking promotion to the next tier; for Disconfirmed claims it names the criteria that failed. Full 36-criterion scorecards appear in \Cref{appendix:audits}.}
\label{tab:verdict-summary}
\begin{tabular}{@{}>{\raggedright\arraybackslash}p{4.35cm}>{\raggedright\arraybackslash}p{2.6cm}>{\raggedright\arraybackslash}p{5.9cm}@{}}
\toprule
\textbf{Claim} & \textbf{Verdict} & \textbf{Gap to next tier} \\
\midrule
Induction heads: in-context token copying \citep{olsson2022context} & Triangulated & Head-evaluator reliability (M1); selection correction over the head pool (M7) \\
\midrule
Copy suppression, L10H7 \citep{mcdougall2023copy} & Mech.\ Supported & Dissociation from backup heads (I6); positive control (M5); top-5\% selection slices (M7) \\
Greater-than circuit \citep{hanna2023greater} & Mech.\ Supported & Second discovery instrument (C3); dissociation (I6); second model (E4) \\
Fourier multiplication in modular addition \citep{nanda2023progress} & Mech.\ Supported & Selection over the frequency pool (M7); positive control (M5) \\
Refusal direction \citep{arditi2024refusal} & Mech.\ Supported & Second estimator for the direction (C3); dissociation (I6); threshold sensitivity (M3) \\
Successor heads \citep{gould2023successor} & Mech.\ Supported & Necessity measured rather than ranked (I1); positive control (M5); selection correction (M7) \\
Superposition in toy models \citep{elhage2022superposition} & Mech.\ Supported & Converse dissociation arm (I6); selection correction (M7); non-synthetic inputs (E2) \\
Global workspace / J-space \citep{gurnee2026verbalizable} & Mech.\ Supported & Converse dissociation (I6); confounder bound (I8); token-set selection (M7) \\
\midrule
Docstring circuit \citep{heimersheim2023docstring} & Causally Suggestive & Specificity on non-docstring behavior (I4); leave-one-out (I3); random head-set baseline (M2) \\
Gender bias circuits \citep{vig2020gender} & Causally Suggestive & Specificity, with no second outcome to register on (I4); necessity inexpressible in the design (I1); selection correction (M7) \\
IOI circuit \citep{wang2023interpretability} & Causally Suggestive & Independent evidence (C3); dissociation (I6); method and prompt robustness (M3, M6) \\
Othello board-state representation \citep{li2023othello} & Causally Suggestive & Necessity by removal (I1); second intervention operator (E1); positive control (M5) \\
SAE features as units of analysis \citep{cunningham2024sparse} & Causally Suggestive & Specificity (I4); independent convergence (C3); calibration of the scoring instrument (M1, M3, M5, M7) \\
\midrule
Probing classifiers \citep{belinkov2022probing} & Proposed & Necessity: the method performs no removal of its own (I1, E1) \\
\midrule
Induction heads: general in-context learning \citep{olsson2022context} & Disconfirmed & Tested and failed: nothing at scale separates induction heads from what forms alongside them (E4, I5); the measure does not separate the construct from few-shot accuracy (C4) \\
Knowledge neurons \citep{dai2022knowledge} & Disconfirmed & Tested and failed: correlation reported, storage claimed (V2); the construct never separates from its neighbor (C4) \\
\bottomrule
\end{tabular}
\end{table}
\endgroup

\subsection{Worked example: induction heads}
\label{sec:example-induction}

\citet{olsson2022context} identify a two-head composition mechanism for in-context copying: previous-token heads attend to the token before a repeated sequence, and induction heads use this signal to predict the next token. The paper advances two claims of different reach, and the evidence has since separated them.

The narrower claim---induction heads implement in-context token copying---reaches Triangulated. Causal ablation now extends to 70B: removing the top 10\% of induction heads in Llama2-70B collapses 8-shot accuracy from 80.57\% to 2.50\%, against 79.13\% for a matched random-head ablation \citep{yang2025unifying}, satisfying E4 by intervention. The construct is defined on repeated sequences of uniformly random tokens, making the claim distribution-free by construction (E2), and \citet{feucht2025dual} later obtained a double dissociation (I6): mean-ablating concept induction heads collapses word-level translation while leaving nonsense copying intact, and the converse holds.

The broader claim---that induction heads are the mechanistic source of \emph{general} in-context learning---is Disconfirmed. At 70--72B the heads carrying abstract in-context reasoning form a disjoint set from induction heads, correlating $r=0.11$ with prefix-matching scores and $r=0.86$ with function-vector scores \citep{webb2025emergent}. Suppressing induction-head formation during training leaves abstract in-context learning intact on 13 of 21 tasks \citep{sahin2025incontext}. And once function-vector heads are preserved, ablating induction heads above 1B is comparable to ablating random heads \citep{yin2025which}.

This is a scope refutation the authors anticipated: their Argument 6 names the condition under which non-induction heads would account for more of in-context learning than induction heads do. No replication was attempted and failed; the narrow mechanism replicates, and the reach asserted for it does not. One paper, two claims, two tiers---the reason the audit scores claims rather than papers.

\subsection{Worked example: the IOI circuit}
\label{sec:example-ioi}

\citet{wang2023interpretability} identify 26 attention heads grouped into seven classes for indirect object identification. Sufficiency is the headline: mean-ablating everything outside the circuit leaves 87\% of the logit difference. Three findings move the claim off Mechanistically Supported, and none of them is a shortage of evidence. First, rival specifications explain the behavior as well: the paper's own na\"ive baseline circuit reaches comparable faithfulness, and \citet{chen2026circuits} exhibit two IOI subgraphs at 100\% task accuracy sharing 4.1\% of their edges. Second, the faithfulness number is a property of the protocol as much as of the circuit---\citet{miller2024transformer} vary six methodological choices and the same quantity ranges from below 0\% to well over 100\%. Third, over $10^6$ benign clean/corrupted pairs from the origin task itself, the model-circuit KL has mean 5.15 and maximum 14.64 \citep{uitdebos2024adversarial}.

Each claim is scored at its most lenient reading, and the lenient reading here is that these 26 heads carry the indirect-object computation. Rival subgraphs are compatible with it, since uniqueness belongs to the stronger reading; method disagreement (E1) caps the claim instead, and no reading escapes that. The localization is real and reproducible, and which subgraph constitutes \emph{the} mechanism is open.

\subsection{Recurring validity gaps}
\label{sec:gaps}

\paragraph{Unattempted criteria (M7, I6, I10).} Within the audit set the systematic absences are near-total. Selection correction (M7) is Untested in all sixteen claims, even though every one of them states its candidate pool---144 heads, 512 neurons, 56 frequencies---and then reports the selected item's score as its property. What the criterion asks for is a headline number recomputed on an unselected slice; it gates no tier, bearing on the size of an effect rather than on whether the mechanism is there. Double dissociation (I6) and rescue reversibility (I10) are Untested in fifteen and thirteen of sixteen respectively, and both have in-field precedent: \citet{feucht2025dual} run the crossed design, and causal tracing performs corrupt-then-restore as its basic operation \citep{meng2022locating}. The instruments exist; the audited claims leave them unrun.

\paragraph{Attempted but unconfirmed criteria (M4, I5, C4, M3, I3).} A second group is attempted almost everywhere and confirmed almost nowhere. Calibration (M4) is attempted in all sixteen claims and Confirmed in none; fourteen land at Partially confirmed, because a metric is reported in interpretable units and never checked against its nearest false positive. Rival mechanism exclusion (I5) is attempted in fifteen and Confirmed in none. Discriminant validity (C4), stability (M3) and minimality (I3) repeat the pattern at eleven to fourteen attempts each. This group is the more tractable of the two, because the experiments are already being run and stop short of the standard. Method-conditionality is the recurring form: the IOI faithfulness range is one instance, and the greater-than circuit's relative faithfulness flips sign under one ablation family. Ablation type is part of the causal claim, not an implementation detail.

\paragraph{No claim reaches Validated.} The Validated tier requires all five validity types to be addressed; no claim in the audit set meets the bar, and induction heads come closest, capped by head-evaluator scores that are single point estimates. The gap is enumerable rather than philosophical: measurement-validity infrastructure (bootstrap stability, seed variance, baseline separation) and interpretive discipline (level declaration, checks on unlicensed labels). A compiled transformer \citep{lindner2023interp} has ground truth by construction; running the full suites on one would produce the first Validated claim and calibrate the framework against a known answer.

% ============================================================
\section{Recommended Practice}
\label{sec:practice}

A validity framework changes nothing if its criteria go unrun, and psychology supplies the measured cost: under the full validation battery it developed, 4\% of examined scales passed, against 88\% under the modal practice of reporting internal consistency alone \citep{hussey2020hidden}; the sixteen audited claims sit in the same position, and most of the unrun tests need no artifact their authors lacked. Five changes to research practice would have the highest impact on claim quality:
\begin{enumerate}[nosep]
\item \emph{Test specificity by default.} Report performance on at least two unrelated tasks alongside the target task; this costs one forward pass per task and addresses the most common gap (I4).
\item \emph{Report multiple ablation methods.} A faithfulness number from one ablation method is a lower bound on the causal claim, not the claim itself.
\item \emph{Calibrate metrics against baselines.} Report the metric score on a random circuit of the same size, an untrained model, and a shuffled-label baseline.
\item \emph{Declare description mode and scope.} This costs zero experiments and prevents interpretive inflation.
\item \emph{Report non-uniqueness.} If alternative circuits achieve comparable faithfulness, report this as a finding; a single reported circuit is a selection from a space the reader cannot see.
\end{enumerate}

We propose three adoption mechanisms: \emph{validity statements}, in which authors declare which tier their claims reach and why; \emph{reviewer checklists}, modeled on CONSORT (Is the description mode declared? Is the claimed tier consistent with the evidence tier?); and \emph{claim audits} in the format of \Cref{appendix:audits}. In the IICARus trial, no manuscript in either arm achieved full compliance with the ARRIVE guidelines \citep{hair2019randomised}, so validity statements and claim audits are proposed as reporting units rather than compliance instruments. For deployment of safety-relevant findings, the tier supplies a minimum standard: deploying a result at Proposed or Causally Suggestive as a monitor is the interpretability equivalent of prescribing a drug that passed Phase~I but not Phase~III.

% ============================================================
\section{Limitations}
\label{sec:limitations}

This work is a theoretical contribution: the verdicts rest on published evidence, not new experiments, and the audited claims predominantly target GPT-2 Small. The sixteen verdicts were assigned after the framework was designed---a retrofitting risk, and the failure mode the framework is designed to detect---so pre-registered application to new claims is planned. Aggregation from criterion-level results to tiers involves judgment, no formal inter-rater reliability study has been conducted, and the tier thresholds are calibrated against the case studies and source-discipline precedents rather than derived from first principles. On calibration of the framework itself, we commit to a ground-truth test: modular addition supplies two rival algorithms known independently of any audit \citep{zhong2023clock}, and the measurement and interpretive suites should certify the algorithm the model implements and deny that tier to the other; a framework that certifies both, or neither, is miscalibrated.

\section{Conclusion}
\label{sec:conclusion}

Mechanistic Validity evaluates mechanistic claims through 36 criteria across five validity types in dependency order. The structured profile makes gaps visible and the path to closing them concrete, and the audit format is proposed as a standard reporting unit.

% ============================================================
\bibliographystyle{plainnat}
\bibliography{references}

\clearpage
\appendix

\section{Description Modes}
\label{appendix:modes}

\begin{table}[H]
\centering
\small
\caption{Seven description modes, ordered from most committal (computational) to least committal (topographic). Implementational-statistical sits partly outside the sequence: it characterizes activation distributions rather than asserting more about a component.}
\label{tab:modes}
\renewcommand{\arraystretch}{1.15}
\begin{tabular}{@{}rl>{\raggedright\arraybackslash}p{4.2cm}>{\raggedright\arraybackslash}p{4.6cm}@{}}
\toprule
& \textbf{Mode} & \textbf{What it asks} & \textbf{Example} \\
\midrule
1 & Computational & What function does the system compute, and why? & ``GPT-2 predicts the indirect object'' \\
2 & Algorithmic & What procedure does it execute? & ``A detect-inhibit-copy algorithm'' \\
3 & Representational & What information is encoded, and in what geometry? & ``The IO name is linearly encoded at layer 8'' \\
4 & Impl.-Functional & What transformation does each component perform? & ``Head 9.9 copies names to the output'' \\
5 & Impl.-Connectomic & How are components wired? & ``S-inhibition feeds name movers via residual stream'' \\
6 & Impl.-Statistical & What do activations look like? & ``Head 9.9 attends strongly to IO position'' \\
7 & Impl.-Topographic & Which components are involved? & ``Heads 9.9, 9.6, 10.0'' \\
\bottomrule
\end{tabular}
\end{table}

\section{Cross-Field Failure Map}
\label{appendix:cross-science}

\Cref{tab:cross-science} maps the validity failures of \Cref{sec:motivation} to the criteria that detect them, alongside their mechanistic-interpretability counterparts.

\begin{table}[H]
\centering
\small
\tighttable
\caption{Validity failures across fields, mapped to Mechanistic Validity criteria. Each failure involves a strong individual result with an untested link in the inference chain.}
\label{tab:cross-science}
\renewcommand{\arraystretch}{1.15}
\begin{tabular}{@{}l>{\raggedright\arraybackslash}p{2.0cm}>{\raggedright\arraybackslash}p{2.5cm}>{\raggedright\arraybackslash}p{5.5cm}@{}}
\toprule
\textbf{Case} & \textbf{Field} & \textbf{Criterion} & \textbf{Untested link} \\
\midrule
Dead salmon fMRI & Neuroscience & M2 Baseline sep. & No null distribution; 70\% false positive rate \\
Cold fusion & Physics & M1/M2 Reliability & Instrument error; \$100M spent \\
BOLD signal & Neuroscience & M1 Reliability & Hemodynamics $\neq$ neural activity \\
SAEBench metrics & MI & M2 Baseline sep. & One metric scores \emph{higher} on a random-init model \\
Nonlinear IIA & MI & M2 Baseline sep. & 100\% IIA on random models \\
\midrule
Cardiac stents & Medicine & I4 Specificity & Sham-controlled: zero benefit \\
Knowledge neurons & MI & I4 Specificity & Fact updating hurts unrelated relations more than random neurons do \\
Candidate genes & Genetics & I7 Confound ctrl & 18 genes $=$ random in $n=620$K \\
COVID-19 ML & ML & I7 Confound ctrl & 0/61 models clinically usable \\
\midrule
IOI circuit & MI & E1 Interv.\ reach & Same quantity spans below 0\% to over 100\% across methods \\
Shortcut learning & ML & E2/C4 Generalization & Texture $\neq$ shape understanding \\
\midrule
Emergent abilities & ML & M4 Calibration & Metric created the phase transition \\
Othello ``world model'' & MI & V4 Interp.\ inflation & Label never contrastively tested against ``causally-used state summary'' \\
\bottomrule
\end{tabular}
\end{table}

\clearpage
\section{Failure Modes and Detecting Criteria}
\label{appendix:failure-modes}

\begingroup
\small
\begin{table}[H]
\centering
\tighttable
\caption{Ten failure modes observed in published mechanistic interpretability research, each mapped to the criteria that detect it.}
\label{tab:failures}
\renewcommand{\arraystretch}{1.15}
\begin{tabular}{@{}l>{\raggedright\arraybackslash}p{3.2cm}>{\raggedright\arraybackslash}p{2.6cm}>{\raggedright\arraybackslash}p{3.9cm}@{}}
\toprule
\textbf{Failure mode} & \textbf{What happens} & \textbf{Criterion} & \textbf{Example} \\
\midrule
Single-metric reporting & Run multiple ablation methods, report only one number & E1 Intervention\newline reach & IOI: one ablation value at origin; the same quantity spans below 0\% to over 100\% across methods \citep{miller2024transformer} \\
\midrule
No baseline control & Report high score without testing a random or untrained model & M2 Baseline\newline separation & SAEBench: Targeted Probe Perturbation scores \emph{higher} on a randomly initialized model \citep{karvonen2025saebench} \\
\midrule
No negative controls & Only test what should pass & I4 Specificity\newline C1 Falsifiability & Docstring circuit: effect on non-docstring behavior never measured \\
\midrule
Untested discriminant validity & Separation from a neighboring construct is presupposed, never measured & C4 Discriminant\newline validity & Gender bias circuits: gender competence is never localized, so no overlap statistic exists \\
\midrule
Off-target effects reported but not read & Off-target numbers are measured, printed, and left uninterpreted & I4 Specificity & Knowledge neurons: fact updating raises inter-relation perplexity 7.2 against 4.3 for random neurons \citep{dai2022knowledge} \\
\midrule
Post-hoc role labeling & Name heads after observed behavior; the label cannot fail & C1 Falsifiability & ``Name mover'' assigned after seeing logit attribution \\
\midrule
No double dissociation & Show necessity but never test the converse & I6 Double\newline dissociation & IOI circuit ablated $\to$ IOI breaks; SVA circuit $\to$ IOI intact never tested \\
\midrule
Interpretive inflation & Label implies intent or algorithm beyond what evidence shows & V4 Unlicensed\newline labeling & ``World model'' for Othello, never contrastively tested against a causally-used state summary \\
\midrule
Scope creep & Claim reach beyond the systems tested & V5 Scope decl.\newline E4 Cross-model & Knowledge neurons: ``can be easily generalized to other pretrained models,'' with no second model run \\
\midrule
Ignoring alternatives & Circuit works, but so does a different set of heads & I5 Rival mech.\newline exclusion & \citet{chen2026circuits} exhibit two IOI subgraphs at 100\% accuracy sharing 4.1\% of their edges \\
\bottomrule
\end{tabular}
\end{table}
\endgroup

\clearpage
\section{Criterion Audits}
\label{appendix:audits}

\subsection*{How to read these audits}

Each claim occupies two pages: a summary page carrying the verdict and how it was reached, and a facing
page carrying the full 36-criterion scorecard. Three conventions govern the scoring.

\textbf{Scope is part of the claim.} A mechanistic claim is not true or false on its own;
it is true or false at a scope. We report the broadest scope at which the claim reaches
each verdict. A claim that holds narrowly and fails broadly is recorded as both, because
reporting only the failure misrepresents the origin paper and reporting only the success
misrepresents how the claim is used. The scorecard scores the primary reading; the readings table records the others.

\textbf{The verdict comes from gates, not from counts.} The status tally at the foot of
each scorecard is descriptive. Promotion between tiers is decided by named criteria, and
the verdict table reports which one blocks it.

\textbf{Untested splits two ways.} \emph{Unattempted} means the experiment was not run.
\emph{Untestable in the setting as run} means it could not have been. Only the first is a
gap the authors could have closed.

\claimheading{Induction Heads: In-Context Token Copying}

\textbf{Source.} \emph{In-context Learning and Induction Heads} \citep{olsson2022context}.

\textbf{Description.} A two-head composition: a previous-token head writes the identity of the token before a repeat into the residual stream, and an induction head reads it to predict what followed last time. Heads are identified by two behavioral scores on sequences of repeated uniform-random tokens, prefix matching and copying, and the mechanism is derived at the weight level for the two-layer attention-only case. Evaluated on: 34 decoder-only models from one to forty layers, with causal ablation confined to the twelve small models, six attention-only and six with MLPs. Description mode: implementational-functional.

\vspace{4pt}
{\footnotesize
\renewcommand{\arraystretch}{1.0}
\begin{longtable}{@{}llll>{\raggedright\arraybackslash}p{4.6cm}@{}}
\caption{\textbf{Induction heads: token copying.} Full 36-criterion audit. Status: \textbf{C} = Confirmed,
\textbf{PC} = Partially confirmed, \textbf{U} = Untested, \textbf{I} = Inconclusive,
\textbf{D} = Disconfirmed, \textbf{N/A} = Not applicable.}\\
\toprule
\textbf{ID} & \textbf{Criterion} & \textbf{Status} & \textbf{Evidence} \\
\endfirsthead
\toprule
\textbf{ID} & \textbf{Criterion} & \textbf{Status} & \textbf{Evidence} \\
\endhead
\midrule
\multicolumn{4}{@{}l}{\textit{Construct Validity}} \\
C1  & Falsifiability            & C   & Two numerical evaluators with a stated failure mode; three predictions \\
C2  & Structural plausibility   & C   & Weight-level: copying is the OV circuit's positive-eigenvalue property \\
C3  & Convergent validity       & C   & Five families: eigenvalues, evaluators, ablation, scaling, training \\
C4  & Discriminant validity     & PC  & Built to separate induction from repeated-token memorization \\
C5  & Nomological validity      & C   & Derived from the mathematical framework rather than fitted to it \\
C6  & Complementation validity  & U   & Two named roles, ablated one at a time and never together \\
\midrule
\multicolumn{4}{@{}l}{\textit{Measurement Validity}} \\
M1  & Reliability               & U   & Unattempted; one training run per model, no interval on any score \\
M2  & Baseline separation       & PC  & Every head ablated at every snapshot, so the effect has a reference \\
M3  & Stability                 & PC  & Token indices perturbed and the score shown robust to the choice \\
M4  & Calibration               & PC  & Evaluators calibrated against weight-level quantities in the appendix \\
M5  & Sensitivity               & PC  & A known-positive exists: the two-layer model with a derived circuit \\
M6  & Invariance                & C   & 34 models, four series, 1--40 layers, 13M--13B, 200 snapshots each \\
M7  & Selection correction      & U   & Unattempted; heads are scored, then the high scorers are ablated \\
\midrule
\multicolumn{4}{@{}l}{\textit{Internal Validity}} \\
I1  & Necessity                 & C   & Almost all in-context learning in small models comes from these heads \\
I2  & Sufficiency               & PC  & No experiment reconstructs copying from induction heads alone \\
I3  & Minimality                & U   & Unattempted; the authors name the obstacle in marginal effects \\
I4  & Specificity               & PC  & Every head is ablated, so off-target effects are measured throughout \\
I5  & Rival mechanism exclusion & PC  & Basic copying heads named and argued away rather than tested \\
I6  & Double dissociation       & C   & Unattempted at origin; run later by Feucht et al.\ and it succeeds \\
I7  & Confound control          & PC  & Three exogenous confounds checked, including scheduled hyperparameters \\
I8  & Confounding sensitivity   & U   & Unattempted; no sensitivity analysis for an unmeasured confounder \\
I9  & Epistatic interaction     & PC  & Composition is the mechanism; the two-head circuit is non-additive \\
I10 & Rescue reversibility      & U   & Unattempted; pattern-preserving ablation already caches the clean run \\
I11 & Onset coupling            & C   & Across 34 models, heads and the loss bump appear at the same point \\
I12 & Offset coupling           & U   & One Pythia-160m head loses the behavior while the algorithm persists \\
\midrule
\multicolumn{4}{@{}l}{\textit{External Validity}} \\
E1  & Intervention reach        & PC  & One primitive at origin: pattern-preserving zero ablation of a head \\
E2  & Prompt generalization     & C   & The defining stimulus is out of distribution by construction \\
E3  & Cross-task generalization & PC  & The same heads perform literal translation, verified against definition \\
E4  & Cross-model generalization& C   & 34 models across four series, plus two external replications \\
E5  & Graded response           & PC  & Both evaluators are continuous; the ablation is not interpolated \\
E6  & Novel prediction          & C   & Argument 2 predicts an architectural change will move the bump \\
\midrule
\multicolumn{4}{@{}l}{\textit{Interpretive Validity}} \\
V1  & Level declaration         & C   & Declared early and explicitly, by behavior on the defining stimulus \\
V2  & Level-evidence match      & C   & Close for the narrow claim; the abstract's narrow sentence is backed \\
V3  & Alternative level         & PC  & One alternative addressed; the authors offer their own caveat \\
V4  & Unlicensed labeling       & PC  & The label imports an epistemic category, and the import is declared \\
V5  & Scope declaration         & C   & The most thoroughly scoped claim in this audit set \\
\midrule
\multicolumn{4}{@{}l}{\textbf{Total:} 14 Confirmed, 15 Partially confirmed, 7 Untested} \\
\multicolumn{4}{@{}l}{\textbf{Verdict: Triangulated}} \\
\bottomrule
\end{longtable}}

\clearpage
\begin{table}[H]
\centering
\footnotesize
\caption{\textbf{Induction heads: token copying.} Readings of the claim, from the version the authors state to the version the
field cites.}
\label{tab:induction-readings}
\renewcommand{\arraystretch}{1.18}
\begin{tabular}{@{}>{\raggedright\arraybackslash}p{3.3cm}l>{\raggedright\arraybackslash}p{4.4cm}@{}}
\toprule
\textbf{Reading} & \textbf{Verdict} & \textbf{Missing} \\
\midrule
\textbf{Primary.} Induction heads implement in-context copying of a token that followed the same context earlier
  & Triangulated
  & Reliability of the head evaluators (M1) and selection correction over the head pool (M7) \\
\addlinespace[2pt]
The mechanism is the source of in-context learning as measured by token-loss difference
  & Mech.\ Supported
  & Necessity above 345M survives a decorrelating control, but the metric is uncalibrated (M4) \\
\addlinespace[2pt]
The mechanism is the source of in-context learning generally
  & Disconfirmed
  & Nothing --- tested and failed; the bridge from the small-model result to the 13B claim is analogy (V2) \\
\addlinespace[2pt]
The two heads compose in the stated order, each performing the role its name asserts
  & Underdetermined
  & Weight-level derivation covers the two-layer case only; role semantics exceed it elsewhere (C2) \\
\bottomrule
\end{tabular}
\end{table}

\begin{table}[H]
\centering
\footnotesize
\caption{\textbf{Induction heads: token copying.} Verdict: \textbf{Triangulated}. A dash means the tier is reached; a criterion at
Inconclusive, Untested or Disconfirmed blocks promotion.}
\label{tab:induction-verdict}
\renewcommand{\arraystretch}{1.18}
\begin{tabular}{@{}l>{\raggedright\arraybackslash}p{3.6cm}>{\raggedright\arraybackslash}p{4.4cm}@{}}
\toprule
\textbf{Tier} & \textbf{Requires} & \textbf{Missing} \\
\midrule
Proposed & C1--C2; one admissible measurement & --- \\
\addlinespace[2pt]
Causally Suggestive & Necessity (I1); baselines (M2) & --- \\
\addlinespace[2pt]
Mechanistically Supported & Sufficiency (I2); reach (E1); specificity (I4) & --- \\
\addlinespace[2pt]
Triangulated & Convergence (C3); replication (E2/E4); dissociation (I6) & --- \\
\addlinespace[2pt]
Validated & Calibration (M1--M7); interpretive audit (V1--V5); reach (E2, E4) & \textbf{M1} unattempted, one training run per model; \textbf{M7} unattempted; \textbf{I3}, \textbf{I8}, \textbf{I10} unattempted \\
\bottomrule
\end{tabular}
\end{table}


\claimheading{Copy Suppression in Negative Attention Heads}

\textbf{Source.} \emph{Copy Suppression: Comprehensively Understanding an Attention Head} \citep{mcdougall2023copy}.

\textbf{Description.} Negative Head L10H7 in GPT-2 Small suppresses tokens that earlier components have already predicted: an early component predicts a token that already appears in context, the head attends back to it, and its OV circuit writes a negative logit for it. The QK and OV halves are read off the weights before any ablation, and a structured ablation deleting every function of the head except those two preserves 76.9\% of its effect on OpenWebText. Description mode: implementational-functional.

\vspace{4pt}
{\footnotesize
\renewcommand{\arraystretch}{1.0}
\begin{longtable}{@{}llll>{\raggedright\arraybackslash}p{4.6cm}@{}}
\caption{\textbf{Copy suppression.} Full 36-criterion audit. Status: \textbf{C} = Confirmed,
\textbf{PC} = Partially confirmed, \textbf{U} = Untested, \textbf{I} = Inconclusive,
\textbf{D} = Disconfirmed, \textbf{N/A} = Not applicable.}\\
\toprule
\textbf{ID} & \textbf{Criterion} & \textbf{Status} & \textbf{Evidence} \\
\endfirsthead
\toprule
\textbf{ID} & \textbf{Criterion} & \textbf{Status} & \textbf{Evidence} \\
\endhead
\midrule
\multicolumn{4}{@{}l}{\textit{Construct Validity}} \\
C1  & Falsifiability            & C   & Three conjunctive conditions, 80\% pass; weight predictions stated first \\
C2  & Structural plausibility   & C   & Read off the weights: 84.70\% OV, 95.72\% QK, before any ablation \\
C3  & Convergent validity       & C   & Four instruments differing in kind agree: weights, coding, ablation, OOD \\
C4  & Discriminant validity     & PC  & Three rivals separated; only the conjunction discriminates, not either half \\
C5  & Nomological validity      & PC  & Unifies negative name movers with anti-induction; formation unexplained \\
C6  & Complementation validity  & U   & OV and QK split within one head; the across-head trans test is absent \\
\midrule
\multicolumn{4}{@{}l}{\textit{Measurement Validity}} \\
M1  & Reliability               & PC  & Spread across data in 100 percentiles; no interval on any estimate \\
M2  & Baseline separation       & PC  & Mean ablation is the denominator; no null constructed for 76.9\% itself \\
M3  & Stability                 & PC  & Perturbations run and disclosed where they fail the idealization \\
M4  & Calibration               & PC  & KL's failure mode named; the effect is small and sign-unstable \\
M5  & Sensitivity               & U   & Every control is a known-negative; no planted mechanism recovered \\
M6  & Invariance                & PC  & All of OpenWebText under three metrics; no domain or frequency split \\
M7  & Selection correction      & U   & Selection disclosed at every step and corrected at none \\
\midrule
\multicolumn{4}{@{}l}{\textit{Internal Validity}} \\
I1  & Necessity                 & PC  & Direct path localized cleanly; the absolute effect is a thousandth of loss \\
I2  & Sufficiency               & PC  & CSPA preserves 76.9\%, but the figure spans 25--95\% across metrics \\
I3  & Minimality                & PC  & Each half mild alone and the conjunction costs; nothing left to prune \\
I4  & Specificity               & PC  & CSPA discriminates L10H7; its components clear 50\% for many heads \\
I5  & Rival mechanism exclusion & PC  & Two rivals addressed; the perpendicular queryside driver left unknown \\
I6  & Double dissociation       & U   & Materials present in Table 2, never assembled into a dissociation \\
I7  & Confound control          & PC  & Tied embeddings and LayerNorm handled; data confounds untouched \\
I8  & Confounding sensitivity   & U   & Unattempted; nothing bounds an unmeasured confounder \\
I9  & Epistatic interaction     & PC  & One interaction quantified, copy suppression on itself; no grid \\
I10 & Rescue reversibility      & U   & Clean activations already in hand; the restore is never run \\
I11 & Onset coupling            & PC  & No onset experiment; named three times as the one to run \\
I12 & Offset coupling           & U   & Pythia checkpoints already in use; no offset observed anywhere \\
\midrule
\multicolumn{4}{@{}l}{\textit{External Validity}} \\
E1  & Intervention reach        & PC  & Four intervention forms, all activation edits on one head \\
E2  & Prompt generalization     & C   & Averaged over the pretraining distribution, not a template set \\
E3  & Cross-task generalization & C   & IOI and random-token scores correlate across three model families \\
E4  & Cross-model generalization& PC  & Phenomenon replicates in three systems; the mechanism does not \\
E5  & Graded response           & PC  & A genuine dose-response at the query; every other intervention binary \\
E6  & Novel prediction          & C   & Anti-induction on random tokens predicted, then confirmed \\
\midrule
\multicolumn{4}{@{}l}{\textit{Interpretive Validity}} \\
V1  & Level declaration         & C   & Three numbered steps mapped to matrices, with the remainder named \\
V2  & Level-evidence match      & PC  & Matched on mechanism and coverage; when and how much it fires is not \\
V3  & Alternative level         & PC  & Appendix K favors semantic copy suppression over the paper's own name \\
V4  & Unlicensed labeling       & C   & Every label redescribes a measured object; no intentional vocabulary \\
V5  & Scope declaration         & C   & Scope declared where each claim is made, not collected at the end \\
\midrule
\multicolumn{4}{@{}l}{\textbf{Total:} 9 Confirmed, 20 Partially confirmed, 7 Untested} \\
\multicolumn{4}{@{}l}{\textbf{Verdict: Mechanistically Supported}} \\
\bottomrule
\end{longtable}}

\clearpage
\begin{table}[H]
\centering
\footnotesize
\caption{\textbf{Copy suppression.} Readings of the claim, from the version the authors state to the version the
field cites.}
\label{tab:copy_suppression-readings}
\renewcommand{\arraystretch}{1.18}
\begin{tabular}{@{}>{\raggedright\arraybackslash}p{3.3cm}l>{\raggedright\arraybackslash}p{4.4cm}@{}}
\toprule
\textbf{Reading} & \textbf{Verdict} & \textbf{Missing} \\
\midrule
\textbf{Primary.} Copy suppression is the main role of Negative Head L10H7 in GPT-2 Small, covering at least 76.9\% of its direct effect
  & Mechanistically Supported
  & A dissociation from the backup-head mechanism the same ablation moves (I6), a known-strength positive control fixing the method's floor (M5), and correction for the top-5\% slices the headline figures are computed on (M7) \\
\addlinespace[2pt]
The token the head suppresses is the source token itself
  & Disconfirmed
  & Nothing --- tested and failed: in 42.00\% of large-attention pairs the source token is suppressed without being the most suppressed, and a semantically related token takes first place in 90\% of those (M6); the authors conclude the mechanism is better read as semantic copy suppression and keep the coarser name elsewhere (V3) \\
\addlinespace[2pt]
Copy suppression is what separates L10H7 from the other heads in its layers
  & Underdetermined
  & Either half of the instrument alone clears 50\% recovered KL for many layer 9--11 heads, so the conjunction rather than the mechanism does the separating (I4, C4); the queryside direction perpendicular to the IO unembedding matters more than the parallel one and is left uncharacterized (I5) \\
\addlinespace[2pt]
Negative Heads do copy suppression across models
  & Proposed
  & The structured ablation is never run outside GPT-2 Small (E4); GPT-2 Medium recovers two of its three most negative heads, Pythia's copy suppression is weaker, and Stanford GPT-2 Small E's analogue attends to IO and S2 equally (M1) \\
\bottomrule
\end{tabular}
\end{table}

\begin{table}[H]
\centering
\footnotesize
\caption{\textbf{Copy suppression.} Verdict: \textbf{Mechanistically Supported}. A dash means the tier is reached; a criterion at
Inconclusive, Untested or Disconfirmed blocks promotion.}
\label{tab:copy_suppression-verdict}
\renewcommand{\arraystretch}{1.18}
\begin{tabular}{@{}l>{\raggedright\arraybackslash}p{3.6cm}>{\raggedright\arraybackslash}p{4.4cm}@{}}
\toprule
\textbf{Tier} & \textbf{Requires} & \textbf{Missing} \\
\midrule
Proposed & C1--C2; one admissible measurement & --- \\
\addlinespace[2pt]
Causally Suggestive & Necessity (I1); baselines (M2) & --- \\
\addlinespace[2pt]
Mechanistically Supported & Sufficiency (I2); reach (E1); specificity (I4) & --- \\
\addlinespace[2pt]
Triangulated & Convergence (C3); replication (E2/E4); dissociation (I6) & \textbf{I6} unattempted; the backup-head mechanism is never shown intact under an ablation that removes copy suppression, and the structured ablation is never applied to it \\
\addlinespace[2pt]
Validated & Calibration (M1--M7); interpretive audit (V1--V5); reach (E2, E4) & \textbf{M5}, \textbf{M7} unattempted, and the behavioral and faithfulness figures are computed on a top-5\% slice chosen by the head's own effect; \textbf{I8}, \textbf{I10}, \textbf{I12} unattempted \\
\bottomrule
\end{tabular}
\end{table}


\claimheading{Greater-Than Circuit in Year-Span Prediction}

\textbf{Source.} \emph{How Does GPT-2 Compute Greater-Than?: Interpreting Mathematical Abilities in a Pre-Trained Language Model} \citep{hanna2023greater}.

\textbf{Description.} Iterative path patching over GPT-2 small recovers a circuit for prompts that supply a two-digit start year and ask for an end year: attention heads carry the start year to the final position, and MLPs 8--11 raise the probability of every year above it. Feeding the circuit real input while the rest of the model receives 01-input recovers 72.7\% of an 81.7\% probability-difference baseline, and the inverse assignment drives the same quantity to $-36.6\%$. Description mode: implementational, with the algorithmic question in the title left open.

\vspace{4pt}
{\footnotesize
\renewcommand{\arraystretch}{1.0}
\begin{longtable}{@{}llll>{\raggedright\arraybackslash}p{4.6cm}@{}}
\caption{\textbf{Greater-than circuit.} Full 36-criterion audit. Status: \textbf{C} = Confirmed,
\textbf{PC} = Partially confirmed, \textbf{U} = Untested, \textbf{I} = Inconclusive,
\textbf{D} = Disconfirmed, \textbf{N/A} = Not applicable.}\\
\toprule
\textbf{ID} & \textbf{Criterion} & \textbf{Status} & \textbf{Evidence} \\
\endfirsthead
\toprule
\textbf{ID} & \textbf{Criterion} & \textbf{Status} & \textbf{Evidence} \\
\endhead
\midrule
\multicolumn{4}{@{}l}{\textit{Construct Validity}} \\
C1  & Falsifiability            & C   & Probability difference moves +81.7\% to -36.6\%, past indifference \\
C2  & Structural plausibility   & PC  & Flow located at three granularities; the derivation fails two causal tests \\
C3  & Convergent validity       & PC  & Logit lens and path patching agree; single-method where they do not overlap \\
C4  & Discriminant validity     & PC  & Three tasks recruit different components; every contrast is one kind \\
C5  & Nomological validity      & PC  & Placed between memorization and generalization, not derived from either \\
C6  & Complementation validity  & U   & Heads and MLPs divided by role, with no joint ablation to test it \\
\midrule
\multicolumn{4}{@{}l}{\textit{Measurement Validity}} \\
M1  & Reliability               & U   & Dispersion reported twice, both pre-circuit; no interval on any result \\
M2  & Baseline separation       & PC  & A size-matched comparison run and reported only in prose \\
M3  & Stability                 & U   & Membership read off heatmaps; no threshold stated, none swept \\
M4  & Calibration               & PC  & Outcome metrics calibrated; the faithfulness quantity is not \\
M5  & Sensitivity               & U   & Two experiments, both on the real model; no planted mechanism \\
M6  & Invariance                & PC  & Per-condition variation visible; one tokenizer, one model, one frame \\
M7  & Selection correction      & U   & Every candidate screened and disclosed; no multiplicity control \\
\midrule
\multicolumn{4}{@{}l}{\textit{Internal Validity}} \\
I1  & Necessity                 & C   & A sign reversal rather than degradation: +81.7\% to -36.6\% \\
I2  & Sufficiency               & C   & Complement recovers 72.7\% under conditions chosen to favor failure \\
I3  & Minimality                & PC  & Most of MLP 10 ablates freely; one result runs the other way \\
I4  & Specificity               & PC  & Fires on less-than prompts its label does not cover \\
I5  & Rival mechanism exclusion & I   & Two rivals tested and both left live, the lookup table included \\
I6  & Double dissociation       & U   & The makings of one half; the converse is never run \\
I7  & Confound control          & PC  & Tokenization confound designed out; frequency handled by restriction \\
I8  & Confounding sensitivity   & U   & Unattempted; nothing bounds an unmeasured confounder \\
I9  & Epistatic interaction     & PC  & Non-additivity is the neuron-level finding; no interaction statistic \\
I10 & Rescue reversibility      & U   & Direction fixed by path patching; the restore is never framed \\
I11 & Onset coupling            & N/A & No published checkpoints, so onset cannot be timed at all \\
I12 & Offset coupling           & N/A & One checkpoint observed once; no interval across which to lapse \\
\midrule
\multicolumn{4}{@{}l}{\textit{External Validity}} \\
E1  & Intervention reach        & PC  & Three intervention families, but one counterfactual behind them all \\
E2  & Prompt generalization     & PC  & Broad inside one frame, and the frame is held fixed by design \\
E3  & Cross-task generalization & PC  & Transfers across the year frame at 67.8--98.8\% and fails outside it \\
E4  & Cross-model generalization& D   & One model throughout, named as a limitation in the conclusion \\
E5  & Graded response           & PC  & Circuit extent is graded; intervention strength has no scale to sweep \\
E6  & Novel prediction          & PC  & Prediction confirmed at 98.8\% and 88.9\%; one corollary falsified \\
\midrule
\multicolumn{4}{@{}l}{\textit{Interpretive Validity}} \\
V1  & Level declaration         & PC  & The declaration is implementational; the title asks an algorithmic question \\
V2  & Level-evidence match      & PC  & Matched on flow and topography, not on the claim the title advertises \\
V3  & Alternative level         & PC  & The lookup-table reading is entertained and bounded, never excluded \\
V4  & Unlicensed labeling       & C   & Every component label redescribes a measurement made first \\
V5  & Scope declaration         & C   & Scope declared before the evidence and again after it \\
\midrule
\multicolumn{4}{@{}l}{\textbf{Total:} 5 Confirmed, 19 Partially confirmed, 1 Inconclusive, 8 Untested, 1 Disconfirmed, 2 Not applicable} \\
\multicolumn{4}{@{}l}{\textbf{Verdict: Mechanistically Supported}} \\
\bottomrule
\end{longtable}}

\clearpage
\begin{table}[H]
\centering
\footnotesize
\caption{\textbf{Greater-than circuit.} Readings of the claim, from the version the authors state to the version the
field cites.}
\label{tab:greater_than-readings}
\renewcommand{\arraystretch}{1.18}
\begin{tabular}{@{}>{\raggedright\arraybackslash}p{3.3cm}l>{\raggedright\arraybackslash}p{4.4cm}@{}}
\toprule
\textbf{Reading} & \textbf{Verdict} & \textbf{Missing} \\
\midrule
\textbf{Primary.} A subgraph of attention heads and MLPs 8--11 performs greater-than on year-span prompts in GPT-2 small
  & Mechanistically Supported
  & A second instrument for circuit discovery (C3), a dissociation from the components serving the tasks the circuit fails (I6), and a second model (E4) \\
\addlinespace[2pt]
The circuit computes the comparison rather than retrieving memorized year associations
  & Underdetermined
  & The paper's own structured-number-representation account fails two causal tests and memorization is never tested (I5); no experiment separates a lookup table from a computation (V3) \\
\addlinespace[2pt]
The circuit implements greater-than as such, firing where the relation applies and not elsewhere
  & Disconfirmed
  & Nothing --- tested and failed: two prompts requiring less-than recruit the same circuit and the model answers greater, and a non-monotone sequence case does the same (I4, E3) \\
\addlinespace[2pt]
The same circuit is the mechanism for greater-than beyond GPT-2 small
  & Disconfirmed
  & Nothing --- tested and failed: the disconfirming evidence is post-origin and in three other model families (E4), and the origin itself runs GPT-2 small only and names the gap \\
\bottomrule
\end{tabular}
\end{table}

\begin{table}[H]
\centering
\footnotesize
\caption{\textbf{Greater-than circuit.} Verdict: \textbf{Mechanistically Supported}. A dash means the tier is reached; a criterion at
Inconclusive, Untested or Disconfirmed blocks promotion.}
\label{tab:greater_than-verdict}
\renewcommand{\arraystretch}{1.18}
\begin{tabular}{@{}l>{\raggedright\arraybackslash}p{3.6cm}>{\raggedright\arraybackslash}p{4.4cm}@{}}
\toprule
\textbf{Tier} & \textbf{Requires} & \textbf{Missing} \\
\midrule
Proposed & C1--C2; one admissible measurement & --- \\
\addlinespace[2pt]
Causally Suggestive & Necessity (I1); baselines (M2) & --- \\
\addlinespace[2pt]
Mechanistically Supported & Sufficiency (I2); reach (E1); specificity (I4) & --- \\
\addlinespace[2pt]
Triangulated & Convergence (C3); replication (E2/E4); dissociation (I6) & \textbf{C3}: every component enters the circuit through one patching method; \textbf{I6} unattempted, with the rival component set located but never intervened on; \textbf{E4} Disconfirmed \\
\addlinespace[2pt]
Validated & Calibration (M1--M7); interpretive audit (V1--V5); reach (E2, E4) & \textbf{M1}, \textbf{M3}, \textbf{M5}, \textbf{M7} unattempted, and the metric that selects the components also scores the circuit they compose; \textbf{I8}, \textbf{I10} unattempted; \textbf{E4} Disconfirmed \\
\bottomrule
\end{tabular}
\end{table}


\claimheading{Fourier Multiplication in Modular Addition}

\textbf{Source.} \emph{Progress Measures for Grokking} \citep{nanda2023progress}.

\textbf{Description.} A one-layer transformer trained on $a+b \bmod 113$ maps each input to $\cos$ and $\sin$ at five key frequencies, multiplies those components so the products encode $a+b$, and reads the sum off at the same five frequencies. The frequencies are identified from the embedding's Fourier spectrum; the account is validated by ablating in Fourier space and by two progress measures, restricted and excluded loss, tracked across training. Restricted and excluded loss move before test accuracy does, which is what makes the account predictive rather than descriptive. Description mode: algorithmic.

\vspace{4pt}
{\footnotesize
\renewcommand{\arraystretch}{1.0}
\begin{longtable}{@{}llll>{\raggedright\arraybackslash}p{4.6cm}@{}}
\caption{\textbf{Fourier multiplication.} Full 36-criterion audit. Status: \textbf{C} = Confirmed,
\textbf{PC} = Partially confirmed, \textbf{U} = Untested, \textbf{I} = Inconclusive,
\textbf{D} = Disconfirmed, \textbf{N/A} = Not applicable.}\\
\toprule
\textbf{ID} & \textbf{Criterion} & \textbf{Status} & \textbf{Evidence} \\
\endfirsthead
\toprule
\textbf{ID} & \textbf{Criterion} & \textbf{Status} & \textbf{Evidence} \\
\endhead
\midrule
\multicolumn{4}{@{}l}{\textit{Construct Validity}} \\
C1  & Falsifiability            & C   & Each key frequency makes a directional ablation prediction that could fail \\
C2  & Structural plausibility   & C   & Weight-level throughout; $W_E$ sparse in the Fourier basis at 6 frequencies \\
C3  & Convergent validity       & PC  & Three families at origin, but all descend from Fourier-space ablation \\
C4  & Discriminant validity     & PC  & Excluded loss rises during circuit formation while train loss stays flat \\
C5  & Nomological validity      & PC  & Placed inside progress measures and phase changes; links are analogical \\
C6  & Complementation validity  & U   & Five frequencies treated as interchangeable, never compared in pairs \\
\midrule
\multicolumn{4}{@{}l}{\textit{Measurement Validity}} \\
M1  & Reliability               & PC  & Five seeds at the mainline configuration, reported in full (Table 3) \\
M2  & Baseline separation       & C   & All 56 frequencies ablated individually; the five separate by orders \\
M3  & Stability                 & PC  & Seed, data fraction, depth and modulus swept; analysis choices are not \\
M4  & Calibration               & PC  & Loss in natural units; the uniform reference is invoked but never computed \\
M5  & Sensitivity               & U   & Unattempted; a known-negative control with no known-positive \\
M6  & Invariance                & C   & Table 5 repeats the mechanism statistics for 27 models, five conditions \\
M7  & Selection correction      & U   & Unattempted, and circular: frequencies chosen from the spectrum tested \\
\midrule
\multicolumn{4}{@{}l}{\textit{Internal Validity}} \\
I1  & Necessity                 & C   & Key-frequency removal gives 6.5--11 in four seeds; nullspace gives 5.27 \\
I2  & Sufficiency               & C   & Non-key Fourier ablation improves loss 70\%; polynomial substitution costs 3\% \\
I3  & Minimality                & PC  & Every key frequency is load-bearing singly; 79 of 512 neurons miss the cutoff \\
I4  & Specificity               & N/A & One task; the model has no off-target behavior to spare \\
I5  & Rival mechanism exclusion & PC  & Memorization excluded; Clock and Pizza are not \\
I6  & Double dissociation       & U   & Unattempted; no second mechanism shown intact under ablation \\
I7  & Confound control          & PC  & Data fraction, modulus, depth, seed and weight decay all varied \\
I8  & Confounding sensitivity   & U   & Unattempted; no sensitivity analysis for an unmeasured confounder \\
I9  & Epistatic interaction     & PC  & Non-additivity demonstrated: single frequencies have sparse effects \\
I10 & Rescue reversibility      & U   & Unattempted; every intervention substitutes rather than restores \\
I11 & Onset coupling            & C   & Restricted and excluded loss move before test accuracy does \\
I12 & Offset coupling           & U   & Unattempted; no test that behavior disappears as the circuit does \\
\midrule
\multicolumn{4}{@{}l}{\textit{External Validity}} \\
E1  & Intervention reach        & C   & Five intervention primitives at four loci, all agreeing \\
E2  & Prompt generalization     & C   & The input space is enumerable and enumerated: all $113^2$ pairs \\
E3  & Cross-task generalization & PC  & Three further tasks run, but only for whether grokking occurs \\
E4  & Cross-model generalization& PC  & One- and two-layer transformers only at origin \\
E5  & Graded response           & U   & Unattempted; ablation is binary per component, with no interpolation \\
E6  & Novel prediction          & C   & The progress measures are derived, then shown to predict the curve \\
\midrule
\multicolumn{4}{@{}l}{\textit{Interpretive Validity}} \\
V1  & Level declaration         & PC  & Level implied by usage rather than declared \\
V2  & Level-evidence match      & PC  & Algorithm-level claim backed at the weight level for most of the model \\
V3  & Alternative level         & PC  & Memorization addressed and refuted; the basis alternative is not \\
V4  & Unlicensed labeling       & C   & Four labels restate measurements; the algorithm name outruns what separates it \\
V5  & Scope declaration         & C   & \S6 states the limits before any generalization is implied \\
\midrule
\multicolumn{4}{@{}l}{\textbf{Total:} 12 Confirmed, 15 Partially confirmed, 8 Untested, 1 Not applicable} \\
\multicolumn{4}{@{}l}{\textbf{Verdict: Triangulated}} \\
\bottomrule
\end{longtable}}

\clearpage
\begin{table}[H]
\centering
\footnotesize
\caption{\textbf{Fourier multiplication.} Readings of the claim, from the version the authors state to the version the
field cites.}
\label{tab:grokking-readings}
\renewcommand{\arraystretch}{1.18}
\begin{tabular}{@{}>{\raggedright\arraybackslash}p{3.3cm}l>{\raggedright\arraybackslash}p{4.4cm}@{}}
\toprule
\textbf{Reading} & \textbf{Verdict} & \textbf{Missing} \\
\midrule
\textbf{Primary.} The network computes $a+b \bmod 113$ by multiplying Fourier components at five key frequencies
  & Mechanistically Supported
  & Selection correction over the frequency pool (M7) and a positive control (M5) \\
\addlinespace[2pt]
The five frequencies are the mechanism at neuron granularity
  & Underdetermined
  & The leave-one-out runs at frequency level; 79 of 512 neurons fail the polynomial fit the account predicts (I3) \\
\addlinespace[2pt]
The account fixes which algorithm the network runs
  & Underdetermined
  & \citet{zhong2023clock} give two algorithms on the same frequencies; the metrics separating them are unrun here (I5) \\
\addlinespace[2pt]
The progress measures explain grokking generally
  & Insufficient
  & Three further tasks are run but only for whether grokking occurs, never for the mechanism (E3) \\
\bottomrule
\end{tabular}
\end{table}

\begin{table}[H]
\centering
\footnotesize
\caption{\textbf{Fourier multiplication.} Verdict: \textbf{Mechanistically Supported}. A dash means the tier is reached; a criterion at
Inconclusive, Untested or Disconfirmed blocks promotion.}
\label{tab:grokking-verdict}
\renewcommand{\arraystretch}{1.18}
\begin{tabular}{@{}l>{\raggedright\arraybackslash}p{3.6cm}>{\raggedright\arraybackslash}p{4.4cm}@{}}
\toprule
\textbf{Tier} & \textbf{Requires} & \textbf{Missing} \\
\midrule
Proposed & C1--C2; one admissible measurement & --- \\
\addlinespace[2pt]
Causally Suggestive & Necessity (I1); baselines (M2) & --- \\
\addlinespace[2pt]
Mechanistically Supported & Sufficiency (I2); reach (E1); specificity (I4) & \textbf{I4} not applicable: one task, no off-target behavior to spare \\
\addlinespace[2pt]
Triangulated & Convergence (C3); replication (E2/E4); dissociation (I6) & \textbf{I6} unattempted; no second mechanism shown intact under key-frequency ablation \\
\addlinespace[2pt]
Validated & Calibration (M1--M7); interpretive audit (V1--V5); reach (E2, E4) & \textbf{M5}, \textbf{M7} unattempted, and the frequency selection is circular; \textbf{I8}, \textbf{I10}, \textbf{E5} unattempted \\
\bottomrule
\end{tabular}
\end{table}


\claimheading{Refusal Mediated by a Single Direction}

\textbf{Source.} \emph{Refusal in Language Models Is Mediated by a Single Direction} \citep{arditi2024refusal}.

\textbf{Description.} A difference in means between activations on harmful and harmless instructions yields one residual-stream direction per model. Projecting that direction out of every matrix that writes to the residual stream stops the model from refusing harmful instructions, and adding it at its extraction layer makes the model refuse harmless ones. Evaluated on: 13 open-weight chat models in 5 families over a 40x parameter range. Description mode: undeclared --- representational in the title, causal-behavioral in the evidence.

\vspace{4pt}
{\footnotesize
\renewcommand{\arraystretch}{1.0}
\begin{longtable}{@{}llll>{\raggedright\arraybackslash}p{4.6cm}@{}}
\caption{\textbf{Refusal direction.} Full 36-criterion audit. Status: \textbf{C} = Confirmed,
\textbf{PC} = Partially confirmed, \textbf{U} = Untested, \textbf{I} = Inconclusive,
\textbf{D} = Disconfirmed, \textbf{N/A} = Not applicable.}\\
\toprule
\textbf{ID} & \textbf{Criterion} & \textbf{Status} & \textbf{Evidence} \\
\endfirsthead
\toprule
\textbf{ID} & \textbf{Criterion} & \textbf{Status} & \textbf{Evidence} \\
\endhead
\midrule
\multicolumn{4}{@{}l}{\textit{Construct Validity}} \\
C1  & Falsifiability            & C   & Two predictions in opposite directions, with the outcome measure fixed first \\
C2  & Structural plausibility   & C   & The edit is enumerated over the five matrix kinds that write to the stream \\
C3  & Convergent validity       & PC  & Two interventions agree on outcome; the estimate itself is never replicated \\
C4  & Discriminant validity     & PC  & The estimator is a harmfulness contrast, and separation is argued by citation \\
C5  & Nomological validity      & PC  & The linear representation hypothesis does not entail the central parameter \\
C6  & Complementation validity  & N/A & One direction, and a single direction has no subdivisions to carve \\
\midrule
\multicolumn{4}{@{}l}{\textit{Measurement Validity}} \\
M1  & Reliability               & PC  & Sampling error over prompts throughout; nothing quantifies the instrument \\
M2  & Baseline separation       & PC  & Five attacks and a fine-tuning comparator, and no baseline for the estimate \\
M3  & Stability                 & U   & Three thresholds and one objective, none of them moved \\
M4  & Calibration               & PC  & Both metrics' failure modes shown by example and never counted \\
M5  & Sensitivity               & PC  & LoRA is an independently established positive and the instrument detects it \\
M6  & Invariance                & C   & 13 models, 5 families, 40x parameters, including where the effect degrades \\
M7  & Selection correction      & U   & Positions crossed with layers, and the search size is never written down \\
\midrule
\multicolumn{4}{@{}l}{\textit{Internal Validity}} \\
I1  & Necessity                 & PC  & 0.95 to 0.01 on all 13 models; unnecessary for refusal produced on instruction \\
I2  & Sufficiency               & PC  & Harmless prompts refused across 13 models, on a direction chosen for sufficiency \\
I3  & Minimality                & PC  & Nothing to prune inside rank one, and no higher rank is compared against \\
I4  & Specificity               & PC  & Six benchmarks and three corpora, with TruthfulQA's fall left unresolved \\
I5  & Rival mechanism exclusion & PC  & Token suppression eliminated on one case; the second rival pre-empted by design \\
I6  & Double dissociation       & U   & One direction and one behavior, tested in both directions of effect \\
I7  & Confound control          & PC  & Disjoint train, validation and evaluation sets; nothing on the extraction axis \\
I8  & Confounding sensitivity   & U   & Unattempted; nothing bounds an unmeasured confounder \\
I9  & Epistatic interaction     & U   & Eight contributors measured one at a time, with no joint test \\
I10 & Rescue reversibility      & U   & The inverse is known in closed form and the restore is never run \\
I11 & Onset coupling            & I   & Base models express the direction as strongly as chat models do \\
I12 & Offset coupling           & U   & Fine-tuning drives refusal to 0.00 and the direction is never re-measured \\
\midrule
\multicolumn{4}{@{}l}{\textit{External Validity}} \\
E1  & Intervention reach        & PC  & A projection and an additive shift, agreeing on outcome and differing on cost \\
E2  & Prompt generalization     & C   & Evaluation prompts disjoint from extraction by construction, across 10 categories \\
E3  & Cross-task generalization & U   & Every behavior tested is refusal or its absence \\
E4  & Cross-model generalization& C   & 13 models, 5 families and a 40x range, not near-copies \\
E5  & Graded response           & U   & Strength is promised in 2.4 and every intervention has a coefficient of 1 \\
E6  & Novel prediction          & PC  & An unrelated attack shows up as suppression, with a matched random control \\
\midrule
\multicolumn{4}{@{}l}{\textit{Interpretive Validity}} \\
V1  & Level declaration         & PC  & Three description modes appear in the first paragraph and none is declared \\
V2  & Level-evidence match      & PC  & The interventions license a lever; the title asserts representational organization \\
V3  & Alternative level         & PC  & The token-suppression reading is retired; the level-of-measure one is untouched \\
V4  & Unlicensed labeling       & PC  & The outcome measure is operationalized; the direction's identity is not \\
V5  & Scope declaration         & C   & Four limitations named, each pointing at what the evidence does not reach \\
\midrule
\multicolumn{4}{@{}l}{\textbf{Total:} 6 Confirmed, 19 Partially confirmed, 1 Inconclusive, 9 Untested, 1 Not applicable} \\
\multicolumn{4}{@{}l}{\textbf{Verdict: Mechanistically Supported}} \\
\bottomrule
\end{longtable}}

\clearpage
\begin{table}[H]
\centering
\footnotesize
\caption{\textbf{Refusal direction.} Readings of the claim, from the version the authors state to the version the
field cites.}
\label{tab:refusal-readings}
\renewcommand{\arraystretch}{1.18}
\begin{tabular}{@{}>{\raggedright\arraybackslash}p{3.3cm}l>{\raggedright\arraybackslash}p{4.4cm}@{}}
\toprule
\textbf{Reading} & \textbf{Verdict} & \textbf{Missing} \\
\midrule
\textbf{Primary.} Ablating one difference-in-means direction stops refusal on harmful instructions and adding it induces refusal on harmless ones, across 13 chat models
  & Mechanistically Supported
  & A second estimator for the direction, since the three implementations are consumers of one difference in means (C3); a dissociation from any other behavior's direction (I6); sensitivity of the extraction to its own three thresholds (M3) \\
\addlinespace[2pt]
The direction carries refusal rather than harmfulness
  & Underdetermined
  & The estimator is the axis along which mean harmful and harmless activations differ, which a harmfulness direction satisfies exactly (V4); the two constructs are separated by citation rather than by experiment (C4); base models that never refuse express the same direction (I11) \\
\addlinespace[2pt]
Refusal is organized along one dimension of the residual stream
  & Underdetermined
  & The interventions license a lever rather than a representational organization, which the authors concede by calling the work an existence proof (V2); no rank above one is ever ablated and the search enumerates single vectors only (I3); refusal is scored as a twelve-substring match, an output property a single direction would control either way (V3) \\
\addlinespace[2pt]
Orthogonalizing the direction removes refusal whatever the prompting condition
  & Disconfirmed
  & Nothing --- tested and failed: with its default system prompt the orthogonalized LLAMA-2 70B yields 4.4\% attack success against 62.9\% without it, so refusal produced on instruction survives an edit that removes the learned propensity (I1, M6) \\
\bottomrule
\end{tabular}
\end{table}

\begin{table}[H]
\centering
\footnotesize
\caption{\textbf{Refusal direction.} Verdict: \textbf{Mechanistically Supported}. A dash means the tier is reached; a criterion at
Inconclusive, Untested or Disconfirmed blocks promotion.}
\label{tab:refusal-verdict}
\renewcommand{\arraystretch}{1.18}
\begin{tabular}{@{}l>{\raggedright\arraybackslash}p{3.6cm}>{\raggedright\arraybackslash}p{4.4cm}@{}}
\toprule
\textbf{Tier} & \textbf{Requires} & \textbf{Missing} \\
\midrule
Proposed & C1--C2; one admissible measurement & --- \\
\addlinespace[2pt]
Causally Suggestive & Necessity (I1); baselines (M2) & --- \\
\addlinespace[2pt]
Mechanistically Supported & Sufficiency (I2); reach (E1); specificity (I4) & --- \\
\addlinespace[2pt]
Triangulated & Convergence (C3); replication (E2/E4); dissociation (I6) & \textbf{C3}: three implementations of two operations, all consuming one unreplicated difference-in-means estimate; \textbf{I6} unattempted, with no second direction ablated to show refusal intact \\
\addlinespace[2pt]
Validated & Calibration (M1--M7); interpretive audit (V1--V5); reach (E2, E4) & \textbf{M3}, \textbf{M7} unattempted, and the selected direction's own selection scores are reported as properties of it; \textbf{I11} Inconclusive, the direction being fully expressed in base models that never refuse; \textbf{I9}, \textbf{I10}, \textbf{E5} unattempted \\
\bottomrule
\end{tabular}
\end{table}


\claimheading{Successor Heads: Ordinal Incrementation}

\textbf{Source.} \emph{Successor Heads: Recurring, Interpretable Attention Heads In The Wild} \citep{gould2023successor}.

\textbf{Description.} One attention head maps an ordinal token to its successor, and the mechanism is a product of four weight matrices --- embedding, $\mathrm{MLP}_0$, the head's OV circuit, unembedding --- evaluated directly, with no activation from any prompt entering it. Mod-10 features recovered from $\mathrm{MLP}_0$ by sparse autoencoders, by a linear probe and by single-neuron ablation carry the index the head increments. Evaluated on: GPT-2, Pythia and Llama-2 from 31M to 12B for behavioral recurrence, with the feature-level case study on Pythia-1.4B L12H0 and two appendix replications. Description mode: implementational-functional.

\vspace{4pt}
{\footnotesize
\renewcommand{\arraystretch}{1.0}
\begin{longtable}{@{}llll>{\raggedright\arraybackslash}p{4.6cm}@{}}
\caption{\textbf{Successor heads.} Full 36-criterion audit. Status: \textbf{C} = Confirmed,
\textbf{PC} = Partially confirmed, \textbf{U} = Untested, \textbf{I} = Inconclusive,
\textbf{D} = Disconfirmed, \textbf{N/A} = Not applicable.}\\
\toprule
\textbf{ID} & \textbf{Criterion} & \textbf{Status} & \textbf{Evidence} \\
\endfirsthead
\toprule
\textbf{ID} & \textbf{Criterion} & \textbf{Status} & \textbf{Evidence} \\
\endhead
\midrule
\multicolumn{4}{@{}l}{\textit{Construct Validity}} \\
C1  & Falsifiability            & C   & A fixed token list and a committed threshold; some models have no head \\
C2  & Structural plausibility   & C   & A product of four matrices tested directly, with no prompt entering it \\
C3  & Convergent validity       & C   & A probe at 0.708 cosine and single neurons, sharing nothing with the SAE \\
C4  & Discriminant validity     & PC  & Top head eight times the runner-up; no neighboring construct attempted \\
C5  & Nomological validity      & PC  & Weak universality imported, tested, and partly returned \\
C6  & Complementation validity  & U   & Eight task classes scored individually and never ablated in pairs \\
\midrule
\multicolumn{4}{@{}l}{\textit{Measurement Validity}} \\
M1  & Reliability               & PC  & Modes over a hundred autoencoders; the central score has no interval \\
M2  & Baseline separation       & PC  & Cutting the head leaves under one percent and returns bigrams instead \\
M3  & Stability                 & PC  & Three perturbations, one aimed at the paper's own premise \\
M4  & Calibration               & PC  & A proportion against a stated threshold; importance gives rank, not magnitude \\
M5  & Sensitivity               & U   & Every control a known-negative; no planted circuit recovered \\
M6  & Invariance                & C   & Three families over three orders of magnitude, failures kept in \\
M7  & Selection correction      & U   & Selection at three levels, disclosed at each and corrected at none \\
\midrule
\multicolumn{4}{@{}l}{\textit{Internal Validity}} \\
I1  & Necessity                 & PC  & Winning head on all 64 prompts, which ranks the head without measuring it \\
I2  & Sufficiency               & C   & The OV circuit alone ranks the successor above every alternative \\
I3  & Minimality                & PC  & Minimal by construction; the droppable component contributes under a percent \\
I4  & Specificity               & PC  & Off-target function measured across five ablation regimes and found present \\
I5  & Rival mechanism exclusion & PC  & Artifact and pre-head computation both excluded by built experiments \\
I6  & Double dissociation       & U   & Neither arm run, and the polysemantic head makes the converse hard \\
I7  & Confound control          & PC  & Tokenization handled by construction; letter collisions explained, not tested \\
I8  & Confounding sensitivity   & U   & Unattempted; nothing bounds an unmeasured confounder \\
I9  & Epistatic interaction     & PC  & Indirect effect bounded under both ablation values; no feature-level grid \\
I10 & Rescue reversibility      & U   & Both ablations reversible and the clean run in hand; restore never run \\
I11 & Onset coupling            & I   & Emergence observed across checkpoints; coupling to behavior is not \\
I12 & Offset coupling           & U   & The checkpoint data exists and is read in the forward direction only \\
\midrule
\multicolumn{4}{@{}l}{\textit{External Validity}} \\
E1  & Intervention reach        & C   & Five intervention forms spanning removal and addition, and they agree \\
E2  & Prompt generalization     & PC  & A constructed token list and 128 natural contexts, the harder one genuinely hard \\
E3  & Cross-task generalization & C   & Eight tasks scored individually; Roman numerals decode though never trained \\
E4  & Cross-model generalization& PC  & Behavior recurs across three families; the mechanism is single-model \\
E5  & Graded response           & PC  & A two-sided dose criterion, and every other ablation all-or-nothing \\
E6  & Novel prediction          & C   & Roman numerals fitted out of every training set and decoded anyway \\
\midrule
\multicolumn{4}{@{}l}{\textit{Interpretive Validity}} \\
V1  & Level declaration         & C   & The level is fixed by the object: a product of four weight matrices \\
V2  & Level-evidence match      & PC  & Weights tested on weights; the recurrence claim rests mostly on behavior \\
V3  & Alternative level         & PC  & Two alternatives closed by experiment; the paper's own third left open \\
V4  & Unlicensed labeling       & C   & Each label defined in the sentence that introduces its measurement \\
V5  & Scope declaration         & PC  & Three limits declared beside the results they bound; one gap remains \\
\midrule
\multicolumn{4}{@{}l}{\textbf{Total:} 10 Confirmed, 18 Partially confirmed, 1 Inconclusive, 7 Untested} \\
\multicolumn{4}{@{}l}{\textbf{Verdict: Mechanistically Supported}} \\
\bottomrule
\end{longtable}}

\clearpage
\begin{table}[H]
\centering
\footnotesize
\caption{\textbf{Successor heads.} Readings of the claim, from the version the authors state to the version the
field cites.}
\label{tab:successor_heads-readings}
\renewcommand{\arraystretch}{1.18}
\begin{tabular}{@{}>{\raggedright\arraybackslash}p{3.3cm}l>{\raggedright\arraybackslash}p{4.4cm}@{}}
\toprule
\textbf{Reading} & \textbf{Verdict} & \textbf{Missing} \\
\midrule
\textbf{Primary.} A single head's effective OV circuit increments an ordinal token to its successor, and mod-10 features carry the index it acts on
  & Mechanistically Supported
  & Necessity measured rather than ranked (I1); a known-positive control (M5); correction for the head, feature and $\lambda$ selections (M7); and the feature account does not reproduce the greater-than bias the same weights show, which the authors state as incomplete (V2) \\
\addlinespace[2pt]
The head is specific to succession
  & Disconfirmed
  & Nothing --- tested and failed: on natural text acronym and greater-than behavior take 23.8\% and 18.9\% of winning cases, and the paper reframes the head as interpretably polysemantic (I4) \\
\addlinespace[2pt]
The mod-10 mechanism recurs across architectures and sizes, as the abstract states
  & Underdetermined
  & Mechanistic recurrence is one case study plus two appendix replications against broad behavioral recurrence (E4); Llama-7B fails the held-out Roman-numeral task where the Pythia models do not (E4, M6), and putting the head on synthesised Roman-numeral representations drops top-1 accuracy to 0.125 (E3, E6) \\
\addlinespace[2pt]
The emergence of successor heads during training explains when the incrementation behavior appears
  & Underdetermined
  & Successor scores are tracked across checkpoints in two model families and the heads are seen to emerge, but nothing is plotted against a behavior or a loss curve, and the phase change the induction-head comparison predicts is absent (I11); no checkpoint interval shows the mechanism lapsing (I12) \\
\bottomrule
\end{tabular}
\end{table}

\begin{table}[H]
\centering
\footnotesize
\caption{\textbf{Successor heads.} Verdict: \textbf{Mechanistically Supported}. A dash means the tier is reached; a criterion at
Inconclusive, Untested or Disconfirmed blocks promotion.}
\label{tab:successor_heads-verdict}
\renewcommand{\arraystretch}{1.18}
\begin{tabular}{@{}l>{\raggedright\arraybackslash}p{3.6cm}>{\raggedright\arraybackslash}p{4.4cm}@{}}
\toprule
\textbf{Tier} & \textbf{Requires} & \textbf{Missing} \\
\midrule
Proposed & C1--C2; one admissible measurement & --- \\
\addlinespace[2pt]
Causally Suggestive & Necessity (I1); baselines (M2) & --- \\
\addlinespace[2pt]
Mechanistically Supported & Sufficiency (I2); reach (E1); specificity (I4) & \textbf{I4} tested and adverse: the head is not specific to succession, with acronym at 23.8\% and greater-than at 18.9\% of winning cases \\
\addlinespace[2pt]
Triangulated & Convergence (C3); replication (E2/E4); dissociation (I6) & \textbf{I6} unattempted, and one head carrying three behaviors makes the converse arm harder to construct; \textbf{E4} mechanistic recurrence is single-model \\
\addlinespace[2pt]
Validated & Calibration (M1--M7); interpretive audit (V1--V5); reach (E2, E4) & \textbf{M5}, \textbf{M7} unattempted, and the case-study head, the mod-10 features and $\lambda$ are each reported at a selected maximum; \textbf{I8}, \textbf{I10}, \textbf{I12} unattempted; \textbf{I11} inconclusive \\
\bottomrule
\end{tabular}
\end{table}


\claimheading{Superposition of Features in Toy Models}

\textbf{Source.} \emph{Toy Models of Superposition} \citep{elhage2022superposition}.

\textbf{Description.} A ReLU network trained to reconstruct sparse features stores more of them than it has dimensions: each feature gets a non-orthogonal direction, the nonlinearity absorbs the interference those directions create, and the linear model differing only in activation function shows no superposition at any sparsity. The tradeoff is mapped over an importance $\times$ sparsity grid against closed-form losses, and the geometries the trained models settle into are uniform polytopes solving a generalized Thomson problem. Evaluated on: toy autoencoders on synthetic features, across four architectures from $(2,1)$ to 400 features in 30 dimensions. Description mode: representational.

\vspace{4pt}
{\footnotesize
\renewcommand{\arraystretch}{1.0}
\begin{longtable}{@{}llll>{\raggedright\arraybackslash}p{4.6cm}@{}}
\caption{\textbf{Superposition.} Full 36-criterion audit. Status: \textbf{C} = Confirmed,
\textbf{PC} = Partially confirmed, \textbf{U} = Untested, \textbf{I} = Inconclusive,
\textbf{D} = Disconfirmed, \textbf{N/A} = Not applicable.}\\
\toprule
\textbf{ID} & \textbf{Criterion} & \textbf{Status} & \textbf{Evidence} \\
\endfirsthead
\toprule
\textbf{ID} & \textbf{Criterion} & \textbf{Status} & \textbf{Evidence} \\
\endhead
\midrule
\multicolumn{4}{@{}l}{\textit{Construct Validity}} \\
C1  & Falsifiability            & C   & Point prediction against a linear control, plus a two-parameter phase diagram \\
C2  & Structural plausibility   & C   & Mechanism derived in closed form, then matched against the trained weights \\
C3  & Convergent validity       & C   & Theory and experiment agree, and three outside groups replicated before publication \\
C4  & Discriminant validity     & PC  & Three-way discrimination against planted features; PCA separated as a limiting case \\
C5  & Nomological validity      & C   & Formal links to compressed sensing, the Thomson problem and distributed codes \\
C6  & Complementation validity  & U   & Ground truth is available by construction and the test is not run \\
\midrule
\multicolumn{4}{@{}l}{\textit{Measurement Validity}} \\
M1  & Reliability               & PC  & Repetition run and reported; aggregation is selection by loss, with no dispersion \\
M2  & Baseline separation       & C   & Linear control cannot superpose, runs at every sparsity, and separates completely \\
M3  & Stability                 & PC  & Sparsity, importance, correlation and size all swept; activation function held fixed \\
M4  & Calibration               & PC  & Two purpose-built measures with stated readings and no external calibration \\
M5  & Sensitivity               & C   & Planted ground-truth features give a known positive at every point \\
M6  & Invariance                & PC  & Holds across sizes and data structures; one activation function and one loss \\
M7  & Selection correction      & U   & Best-of-1000 and worst-discarded selection documented, with no correction applied \\
\midrule
\multicolumn{4}{@{}l}{\textit{Internal Validity}} \\
I1  & Necessity                 & C   & Output nonlinearity, sparsity and negative bias each shown to be required \\
I2  & Sufficiency               & C   & Built from the hypothesized ingredients alone, and superposition appears \\
I3  & Minimality                & PC  & Load-bearing parts named and the smallest case solved; data assumptions untested \\
I4  & Specificity               & N/A & One task and one loss, and the paper says cross-superposition loss comparison fails \\
I5  & Rival mechanism exclusion & PC  & PCA named and characterized as a limiting case; incidental polysemanticity untested \\
I6  & Double dissociation       & U   & Status recorded without a supporting quote from the source \\
I7  & Confound control          & PC  & Data confounds set by construction; optimization left as an uncontrolled one \\
I8  & Confounding sensitivity   & U   & Status recorded without a supporting quote from the source \\
I9  & Epistatic interaction     & C   & Feature interactions are the subject once uniformity is dropped, and each is measured \\
I10 & Rescue reversibility      & U   & Status recorded without a supporting quote from the source \\
I11 & Onset coupling            & PC  & Onset tracked across training; feature dimensionality jumps as the loss drops \\
I12 & Offset coupling           & U   & Offset never run; the learning-dynamics section is limited by the authors' own account \\
\midrule
\multicolumn{4}{@{}l}{\textit{External Validity}} \\
E1  & Intervention reach        & PC  & Four intervention forms including two that remove superposition, cost reported \\
E2  & Prompt generalization     & PC  & Input distribution swept along every axis the theory names, all synthetic \\
E3  & Cross-task generalization & PC  & Two tasks, with computation carrying the weight; representativeness left open \\
E4  & Cross-model generalization& PC  & Four architectures inside the toy family; real-model support is consistency only \\
E5  & Graded response           & C   & Dose-response is the core object: sparsity swept, response monotone and located \\
E6  & Novel prediction          & C   & Two predictions that could have failed: polytope geometry and adversarial fragility \\
\midrule
\multicolumn{4}{@{}l}{\textit{Interpretive Validity}} \\
V1  & Level declaration         & C   & Levels separated before any experiment; the chosen definition flagged as circular \\
V2  & Level-evidence match      & C   & Claims stay at the toy-model level, and every step up from it is marked \\
V3  & Alternative level         & PC  & PCA examined as a limiting case; the optimization alternative raised and left open \\
V4  & Unlicensed labeling       & C   & Metaphors are physical and cashed out; agentive verbs appear but stay scare-quoted \\
V5  & Scope declaration         & C   & Scope declared in the title, in Key Results, in the Discussion and in Open Questions \\
\midrule
\multicolumn{4}{@{}l}{\textbf{Total:} 15 Confirmed, 14 Partially confirmed, 6 Untested, 1 Not applicable} \\
\multicolumn{4}{@{}l}{\textbf{Verdict: Mechanistically Supported}} \\
\bottomrule
\end{longtable}}

\clearpage
\begin{table}[H]
\centering
\footnotesize
\caption{\textbf{Superposition.} Readings of the claim, from the version the authors state to the version the
field cites.}
\label{tab:superposition-readings}
\renewcommand{\arraystretch}{1.18}
\begin{tabular}{@{}>{\raggedright\arraybackslash}p{3.3cm}l>{\raggedright\arraybackslash}p{4.4cm}@{}}
\toprule
\textbf{Reading} & \textbf{Verdict} & \textbf{Missing} \\
\midrule
\textbf{Primary.} A ReLU network on sparse features represents more features than it has dimensions, in geometries set by sparsity and importance
  & Mechanistically Supported
  & A converse arm, with some property shown intact while superposition is removed (I6); correction for the lowest-loss selection the reported geometries are drawn from (M7); an input distribution outside the synthetic family (E2) \\
\addlinespace[2pt]
The reported geometries are properties of the optimum rather than of the optimizer
  & Underdetermined
  & The $m=2$ case is reported as much harder for gradient descent and its solutions are selected by loss (I7, M7); an invited comment in the same article shows the global minimum can have a smaller basin of attraction than nearby local minima, and the paper leaves it open (V3) \\
\addlinespace[2pt]
The phase diagram is fixed by sparsity and importance alone
  & Disconfirmed
  & Nothing --- tested and failed: the paper holds the activation function fixed across every one of its own experiments, and a replication carried in the same article finds the phase diagrams look quite different under other activation functions (M3, M6) \\
\addlinespace[2pt]
Language models represent features in superposition
  & Insufficient
  & Every measurement is on a toy model, and the paper labels its real-model section validation by consistency with existing reports rather than measurement taken here (E4); no natural-data distribution appears anywhere, and Open Questions asks whether real importance and sparsity curves can be estimated at all (E2) \\
\bottomrule
\end{tabular}
\end{table}

\begin{table}[H]
\centering
\footnotesize
\caption{\textbf{Superposition.} Verdict: \textbf{Mechanistically Supported}. A dash means the tier is reached; a criterion at
Inconclusive, Untested or Disconfirmed blocks promotion.}
\label{tab:superposition-verdict}
\renewcommand{\arraystretch}{1.18}
\begin{tabular}{@{}l>{\raggedright\arraybackslash}p{3.6cm}>{\raggedright\arraybackslash}p{4.4cm}@{}}
\toprule
\textbf{Tier} & \textbf{Requires} & \textbf{Missing} \\
\midrule
Proposed & C1--C2; one admissible measurement & --- \\
\addlinespace[2pt]
Causally Suggestive & Necessity (I1); baselines (M2) & --- \\
\addlinespace[2pt]
Mechanistically Supported & Sufficiency (I2); reach (E1); specificity (I4) & \textbf{I4} not applicable: one objective, and superposition is induced by changing the task, so losses at different amounts of it are measured on different tasks \\
\addlinespace[2pt]
Triangulated & Convergence (C3); replication (E2/E4); dissociation (I6) & \textbf{I6} unattempted; no second property is shown intact while superposition is removed, and no manipulation removes another property while leaving superposition in place \\
\addlinespace[2pt]
Validated & Calibration (M1--M7); interpretive audit (V1--V5); reach (E2, E4) & \textbf{M7} unattempted, and the neuron-level figures are the best of a thousand fits while the phase diagram discards the worst of ten; \textbf{I8}, \textbf{I10}, \textbf{I12} unattempted \\
\bottomrule
\end{tabular}
\end{table}


\claimheading{Global Workspace / J-space}

\textbf{Source.} \emph{Verbalizable Representations Form a Global Workspace} \citep{gurnee2026verbalizable}.

\textbf{Description.} A low-dimensional subspace of the residual stream, recovered as the expected Jacobian $J_\ell$ of the final-layer residual with respect to layer $\ell$, is proposed as the set of contents a model can verbally report and reason over. Five properties are enumerated in advance from global workspace theory --- verbal report, directed modulation, internal reasoning, flexible generalization, capacity limits --- and one is used to find the object rather than to score it afterwards. Evaluated on: four closed-weight Claude models with sizes undisclosed, plus base and model-organism checkpoints. Description mode: implementational-functional.

\vspace{4pt}
{\footnotesize
\renewcommand{\arraystretch}{1.0}
\begin{longtable}{@{}llll>{\raggedright\arraybackslash}p{4.6cm}@{}}
\caption{\textbf{Global workspace / J-space.} Full 36-criterion audit. Status: \textbf{C} = Confirmed,
\textbf{PC} = Partially confirmed, \textbf{U} = Untested, \textbf{I} = Inconclusive,
\textbf{D} = Disconfirmed, \textbf{N/A} = Not applicable.}\\
\toprule
\textbf{ID} & \textbf{Criterion} & \textbf{Status} & \textbf{Evidence} \\
\endfirsthead
\toprule
\textbf{ID} & \textbf{Criterion} & \textbf{Status} & \textbf{Evidence} \\
\endhead
\midrule
\multicolumn{4}{@{}l}{\textit{Construct Validity}} \\
C1  & Falsifiability            & C   & Five properties operationalized before any result is reported \\
C2  & Structural plausibility   & C   & $J_\ell$ derived from the architecture, not posited \\
C3  & Convergent validity       & PC  & Three cross-checks, but almost every intervention edits along J-lens \\
C4  & Discriminant validity     & PC  & Separated from ``any direction encoding this concept'' by design \\
C5  & Nomological validity      & C   & Stated inside a named theory with disanalogies enumerated \\
C6  & Complementation validity  & N/A & One subspace, undivided, so the question does not attach \\
\midrule
\multicolumn{4}{@{}l}{\textit{Measurement Validity}} \\
M1  & Reliability               & PC  & Wilson intervals and per-prompt dots throughout; no seed variance \\
M2  & Baseline separation       & C   & The densest negative controls in the set: logit lens and tuned lens \\
M3  & Stability                 & PC  & Method variants swept and robust; corpus size from 1 to 1000 \\
M4  & Calibration               & PC  & Ablations normalized to unablated Sonnet with a smaller-model floor \\
M5  & Sensitivity               & C   & A known-positive control is run and named as one \\
M6  & Invariance                & PC  & Four models, base and organism checkpoints, two further lenses \\
M7  & Selection correction      & U   & Unattempted, and one token set is chosen self-referentially \\
\midrule
\multicolumn{4}{@{}l}{\textit{Internal Validity}} \\
I1  & Necessity                 & C   & Four necessity demonstrations, three with controls \\
I2  & Sufficiency               & PC  & Swaps install counterfactual content at 54--70\% across three models \\
I3  & Minimality                & PC  & Occupancy estimated from where marginal improvement flattens \\
I4  & Specificity               & C   & Tested three ways; a fourteen-task battery maps where it is not needed \\
I5  & Rival mechanism exclusion & PC  & The concept-vector rival is named twice and tested twice \\
I6  & Double dissociation       & U   & Single dissociation run cleanly and repeatedly; no converse \\
I7  & Confound control          & C   & Five designed controls, each aimed at its experiment's alternative \\
I8  & Confounding sensitivity   & U   & Unattempted; specific alternatives answered, no unmeasured-confounder bound \\
I9  & Epistatic interaction     & PC  & Token-level exclusion at 0.09 vs 0.29; two held concepts co-occupy at chance \\
I10 & Rescue reversibility      & PC  & Clamp-to-clean restoration run twice as a mediation control \\
I11 & Onset coupling            & PC  & Present and load-bearing in a base model before any RLHF \\
I12 & Offset coupling           & PC  & Deception signal amplified by coding RL, then attenuated by safety training \\
\midrule
\multicolumn{4}{@{}l}{\textit{External Validity}} \\
E1  & Intervention reach        & PC  & Many intervention forms, nearly all editing along the same directions \\
E2  & Prompt generalization     & C   & Five public benchmarks plus six prompt distributions \\
E3  & Cross-task generalization & C   & A fourteen-task battery maps where the mechanism is required \\
E4  & Cross-model generalization& PC  & Four closed-weight models, one lab, sizes undisclosed; partly replicated externally \\
E5  & Graded response           & C   & Six dose-response designs, including injection strength against rank \\
E6  & Novel prediction          & C   & Five properties enumerated from theory, one used to find the object \\
\midrule
\multicolumn{4}{@{}l}{\textit{Interpretive Validity}} \\
V1  & Level declaration         & C   & Declared before evidence, and the level is named explicitly \\
V2  & Level-evidence match      & PC  & Three levels in play and two matched; the third is asserted \\
V3  & Alternative level         & PC  & Two alternatives addressed and refuted; the workspace framing is not \\
V4  & Unlicensed labeling       & PC  & Disanalogies enumerated and consciousness not adjudicated; hedged \\
V5  & Scope declaration         & C   & Eight limitations named in a dedicated section \\
\midrule
\multicolumn{4}{@{}l}{\textbf{Total:} 14 Confirmed, 18 Partially confirmed, 3 Untested, 1 Not applicable} \\
\multicolumn{4}{@{}l}{\textbf{Verdict: Mechanistically Supported}} \\
\bottomrule
\end{longtable}}

\clearpage
\begin{table}[H]
\centering
\footnotesize
\caption{\textbf{Global workspace / J-space.} Readings of the claim, from the version the authors state to the version the
field cites.}
\label{tab:workspace-readings}
\renewcommand{\arraystretch}{1.18}
\begin{tabular}{@{}>{\raggedright\arraybackslash}p{3.3cm}l>{\raggedright\arraybackslash}p{4.4cm}@{}}
\toprule
\textbf{Reading} & \textbf{Verdict} & \textbf{Missing} \\
\midrule
\textbf{Primary.} A low-dimensional Jacobian subspace mediates which contents the model reports and reasons over
  & Mechanistically Supported
  & A converse dissociation (I6), an unmeasured-confounder bound (I8), and correction for how the token sets were chosen (M7) \\
\addlinespace[2pt]
Any direction encoding the concept would serve equally
  & Disconfirmed
  & Nothing --- named twice and tested twice, and concept vectors do not reproduce the effect \\
\addlinespace[2pt]
The subspace is a global workspace in the sense the theory intends
  & Underdetermined
  & The architectural disanalogies are enumerated by the authors; no experiment separates a workspace from a bottleneck that behaves like one \\
\addlinespace[2pt]
The result holds beyond one model family
  & Insufficient
  & Four models, one laboratory, all closed-weight, effects scale-graded from 54\% to 70\% (E4) \\
\bottomrule
\end{tabular}
\end{table}

\begin{table}[H]
\centering
\footnotesize
\caption{\textbf{Global workspace / J-space.} Verdict: \textbf{Mechanistically Supported}. A dash means the tier is reached; a criterion at
Inconclusive, Untested or Disconfirmed blocks promotion.}
\label{tab:workspace-verdict}
\renewcommand{\arraystretch}{1.18}
\begin{tabular}{@{}l>{\raggedright\arraybackslash}p{3.6cm}>{\raggedright\arraybackslash}p{4.4cm}@{}}
\toprule
\textbf{Tier} & \textbf{Requires} & \textbf{Missing} \\
\midrule
Proposed & C1--C2; one admissible measurement & --- \\
\addlinespace[2pt]
Causally Suggestive & Necessity (I1); baselines (M2) & --- \\
\addlinespace[2pt]
Mechanistically Supported & Sufficiency (I2); reach (E1); specificity (I4) & --- \\
\addlinespace[2pt]
Triangulated & Convergence (C3); replication (E2/E4); dissociation (I6) & \textbf{C3}: nearly every intervention edits along the same directions; \textbf{I6} single rather than double; \textbf{E4} one model family \\
\addlinespace[2pt]
Validated & Calibration (M1--M7); interpretive audit (V1--V5); reach (E2, E4) & \textbf{M7} unattempted, and one token set is chosen self-referentially; \textbf{I8} unattempted \\
\bottomrule
\end{tabular}
\end{table}


\claimheading{Docstring Circuit}

\textbf{Source.} \emph{A circuit for Python docstrings in a 4-layer attention-only transformer} \citep{heimersheim2023docstring}.

\textbf{Description.} Eight attention heads in a 4-layer attention-only transformer compose across three levels to predict the next argument name in a Python docstring: fuzzy previous-token heads and a positional head set up an induction step, and argument movers carry the name from the definition line to the output position. Resample-ablating everything outside the eight leaves 42\% success against the full model's 56\%. Evaluated on: one released toy model, 50 prompts per experiment. Description mode: implementational-functional.

\vspace{4pt}
{\footnotesize
\renewcommand{\arraystretch}{1.0}
\begin{longtable}{@{}llll>{\raggedright\arraybackslash}p{4.6cm}@{}}
\caption{\textbf{Docstring circuit.} Full 36-criterion audit. Status: \textbf{C} = Confirmed,
\textbf{PC} = Partially confirmed, \textbf{U} = Untested, \textbf{I} = Inconclusive,
\textbf{D} = Disconfirmed, \textbf{N/A} = Not applicable.}\\
\toprule
\textbf{ID} & \textbf{Criterion} & \textbf{Status} & \textbf{Evidence} \\
\endfirsthead
\toprule
\textbf{ID} & \textbf{Criterion} & \textbf{Status} & \textbf{Evidence} \\
\endhead
\midrule
\multicolumn{4}{@{}l}{\textit{Construct Validity}} \\
C1  & Falsifiability            & C   & Each head assignment is a signed prediction that could have come out flat \\
C2  & Structural plausibility   & PC  & Argued at the activation level: patterns locate, patching confirms \\
C3  & Convergent validity       & PC  & One instrument carries it; patching and composition share a primitive \\
C4  & Discriminant validity     & U   & The circuit is never compared with one for a neighboring task \\
C5  & Nomological validity      & PC  & Placed inside an existing theory of head types, which does work \\
C6  & Complementation validity  & U   & Three named levels, each ablated alone and never in combination \\
\midrule
\multicolumn{4}{@{}l}{\textit{Measurement Validity}} \\
M1  & Reliability               & U   & Fifty prompts behind every plot; dispersion is never stated \\
M2  & Baseline separation       & PC  & 42\% against a 17\% chance level, with two further reference points \\
M3  & Stability                 & U   & Head membership is varied upward only; nothing re-measures the claim \\
M4  & Calibration               & PC  & Logit difference against the highest wrong answer, recomputed throughout \\
M5  & Sensitivity               & U   & Nothing with a known answer is fed in to see whether it comes back \\
M6  & Invariance                & U   & No conclusion is re-derived under a changed condition and shown to hold \\
M7  & Selection correction      & U   & Components are chosen from the strongest cells of a patching grid \\
\midrule
\multicolumn{4}{@{}l}{\textit{Internal Validity}} \\
I1  & Necessity                 & PC  & Every head is knocked out by resampling and the tests bite \\
I2  & Sufficiency               & PC  & Everything outside the eight heads is ablated and reported straight \\
I3  & Minimality                & PC  & Heads are added rather than removed; one member is undercut anyway \\
I4  & Specificity               & U   & Measured on one behavior; head 1.4 alone is run off-distribution \\
I5  & Rival mechanism exclusion & PC  & Rivals are named first and the prompt set disables two of them \\
I6  & Double dissociation       & U   & No second capability is shown intact while the docstring one fails \\
I7  & Confound control          & PC  & Line counting and repetition are designed out rather than argued away \\
I8  & Confounding sensitivity   & U   & Nothing asks how strong an unmeasured confounder would need to be \\
I9  & Epistatic interaction     & PC  & The composition score is a two-head quantity, so interaction is the subject \\
I10 & Rescue reversibility      & U   & The restoring direction was run and not shown, by the author's own account \\
I11 & Onset coupling            & U   & One pre-trained model off the shelf; no training trajectory in evidence \\
I12 & Offset coupling           & U   & Needs the training sequence that onset coupling also lacks \\
\midrule
\multicolumn{4}{@{}l}{\textit{External Validity}} \\
E1  & Intervention reach        & PC  & Four do-operators beyond resampling, and they agree with it \\
E2  & Prompt generalization     & D   & A subclass of in-distribution inputs makes the circuit diverge from the model \\
E3  & Cross-task generalization & U   & Whether the components do anything for another task is never asked \\
E4  & Cross-model generalization& U   & Transfer is named as future work and no result is reported \\
E5  & Graded response           & PC  & Circuit size tracks performance monotonically across two axes \\
E6  & Novel prediction          & C   & A mechanism proposed for one reason predicts something else, and it holds \\
\midrule
\multicolumn{4}{@{}l}{\textit{Interpretive Validity}} \\
V1  & Level declaration         & C   & A circuit is defined as a subset of components before any result \\
V2  & Level-evidence match      & PC  & Matches the lower declared level and thins at the upper one \\
V3  & Alternative level         & PC  & Line counting is raised as the simpler description, then accepted \\
V4  & Unlicensed labeling       & C   & Each label is introduced where its operation has just been measured \\
V5  & Scope declaration         & C   & The limit is declared before the argument, not conceded after it \\
\midrule
\multicolumn{4}{@{}l}{\textbf{Total:} 5 Confirmed, 15 Partially confirmed, 15 Untested, 1 Disconfirmed} \\
\multicolumn{4}{@{}l}{\textbf{Verdict: Causally Suggestive}} \\
\bottomrule
\end{longtable}}

\clearpage
\begin{table}[H]
\centering
\footnotesize
\caption{\textbf{Docstring circuit.} Readings of the claim, from the version the authors state to the version the
field cites.}
\label{tab:docstring-readings}
\renewcommand{\arraystretch}{1.18}
\begin{tabular}{@{}>{\raggedright\arraybackslash}p{3.3cm}l>{\raggedright\arraybackslash}p{4.4cm}@{}}
\toprule
\textbf{Reading} & \textbf{Verdict} & \textbf{Missing} \\
\midrule
\textbf{Primary.} Eight named heads compose to predict the next argument name in a Python docstring in this 4-layer model
  & Causally Suggestive
  & Specificity on any non-docstring behavior (I4), a leave-one-out over the eight heads (I3), a matched random head set to show the 42\% belongs to these heads (M2), and dispersion on any reported number (M1) \\
\addlinespace[2pt]
The circuit is the model's implementation of Docstring Induction, an algorithm
  & Underdetermined
  & Apportionment between induction and the line-number algorithm, which the authors judge co-implemented and suppress by prompt design rather than refute (I5, V3); the algorithm is inferred from attention patterns rather than tested against its rivals (V2) \\
\addlinespace[2pt]
The circuit behaves like the model on the task it was found on
  & Disconfirmed
  & Nothing --- tested and failed (E2), and on post-origin evidence alone; the post's own remark that Google-style docstrings give similar performance points the other way and carries no number (M3) \\
\addlinespace[2pt]
These heads, or the head types they instantiate, serve other tasks or recur in other models
  & Insufficient
  & Cross-task reuse is never asked (E3), cross-model transfer is named as future work with no model, head set or number reported (E4), and the circuit is never compared with one for a neighboring task (C4) \\
\bottomrule
\end{tabular}
\end{table}

\begin{table}[H]
\centering
\footnotesize
\caption{\textbf{Docstring circuit.} Verdict: \textbf{Causally Suggestive}. A dash means the tier is reached; a criterion at
Inconclusive, Untested or Disconfirmed blocks promotion.}
\label{tab:docstring-verdict}
\renewcommand{\arraystretch}{1.18}
\begin{tabular}{@{}l>{\raggedright\arraybackslash}p{3.6cm}>{\raggedright\arraybackslash}p{4.4cm}@{}}
\toprule
\textbf{Tier} & \textbf{Requires} & \textbf{Missing} \\
\midrule
Proposed & C1--C2; one admissible measurement & --- \\
\addlinespace[2pt]
Causally Suggestive & Necessity (I1); baselines (M2) & --- \\
\addlinespace[2pt]
Mechanistically Supported & Sufficiency (I2); reach (E1); specificity (I4) & \textbf{I4} unattempted: no non-docstring behavior is measured, in a model that also handles natural language; \textbf{I2} recovers 42\% against the full model's 56\% \\
\addlinespace[2pt]
Triangulated & Convergence (C3); replication (E2/E4); dissociation (I6) & \textbf{E2} Disconfirmed; \textbf{E4} unattempted and declared unfinished; \textbf{I6} unattempted; \textbf{C3}: patching and the composition score share one primitive \\
\addlinespace[2pt]
Validated & Calibration (M1--M7); interpretive audit (V1--V5); reach (E2, E4) & \textbf{M1}, \textbf{M3}, \textbf{M5}, \textbf{M6}, \textbf{M7} unattempted --- fourteen of the 35 criteria are Untested; \textbf{E2} Disconfirmed \\
\bottomrule
\end{tabular}
\end{table}


\claimheading{Gender Bias Circuits}

\textbf{Source.} \emph{Causal Mediation Analysis for Interpreting Neural NLP: The Case of Gender Bias} \citep{vig2020gender}.

\textbf{Description.} Causal mediation analysis treats a gender edit to the input as a treatment and each neuron, layer, head or attention weight as a mediator, splitting the effect on a two-pronoun probability ratio into natural direct and indirect effects. The effect is sparse: ten of 144 heads reproduce the effect of intervening on all of them. A randomly initialized GPT-2 Small is carried through as a matched negative control. Evaluated on: five GPT-2 sizes on 17 templates over 169 professions plus WinoBias and Winogender. Description mode: implementational-topographic.

\vspace{4pt}
{\footnotesize
\renewcommand{\arraystretch}{1.0}
\begin{longtable}{@{}llll>{\raggedright\arraybackslash}p{4.6cm}@{}}
\caption{\textbf{Gender bias circuits.} Full 36-criterion audit. Status: \textbf{C} = Confirmed,
\textbf{PC} = Partially confirmed, \textbf{U} = Untested, \textbf{I} = Inconclusive,
\textbf{D} = Disconfirmed, \textbf{N/A} = Not applicable.}\\
\toprule
\textbf{ID} & \textbf{Criterion} & \textbf{Status} & \textbf{Evidence} \\
\endfirsthead
\toprule
\textbf{ID} & \textbf{Criterion} & \textbf{Status} & \textbf{Evidence} \\
\endhead
\midrule
\multicolumn{4}{@{}l}{\textit{Construct Validity}} \\
C1  & Falsifiability            & C   & The answer could have been 100 heads and was 10; one negative case fires \\
C2  & Structural plausibility   & PC  & Mediators taken from the architecture; no weight-level mechanism anywhere \\
C3  & Convergent validity       & C   & Two mediator families, three datasets, four scales -- all one instrument \\
C4  & Discriminant validity     & U   & No second construct is ever localized, so nothing can separate from it \\
C5  & Nomological validity      & C   & Pearl's estimands used under his definitions, with the conditions stated \\
C6  & Complementation validity  & N/A & Mediators carry an index and an estimand, not a labeled role \\
\midrule
\multicolumn{4}{@{}l}{\textit{Measurement Validity}} \\
M1  & Reliability               & PC  & Dispersion reported where it is incidental and absent where it is load-bearing \\
M2  & Baseline separation       & C   & An untrained model matched on everything but training, run at both levels \\
M3  & Stability                 & PC  & Selection rule and outcome scale both varied; the target word is not \\
M4  & Calibration               & PC  & A fixed point at 1 and a sign convention; the effect scale has no referent \\
M5  & Sensitivity               & PC  & Effects track published occupational statistics, but at the wrong level \\
M6  & Invariance                & C   & Six axes tested, and where the invariance stops is reported too \\
M7  & Selection correction      & U   & Pool, rule and shortcut all disclosed; nothing corrected for multiplicity \\
\midrule
\multicolumn{4}{@{}l}{\textit{Internal Validity}} \\
I1  & Necessity                 & N/A & Neither estimand can be put in the form necessity requires \\
I2  & Sufficiency               & PC  & Ten heads of 144 reproduce the all-mediator effect; no capability measured \\
I3  & Minimality                & PC  & Greedy admission by marginal gain; no ablation of the selected ten \\
I4  & Specificity               & U   & A two-element outcome set leaves off-target damage nowhere to register \\
I5  & Rival mechanism exclusion & PC  & The untrained control excludes architecture; the stimulus rival stands \\
I6  & Double dissociation       & U   & One behavior by construction, so the second dissociation cannot be run \\
I7  & Confound control          & PC  & Definitional gender dropped by crowdsourced rating; frequency untouched \\
I8  & Confounding sensitivity   & U   & Unattempted; nothing bounds an unmeasured confounder \\
I9  & Epistatic interaction     & C   & Concurrent against summed intervention, run at both levels, and it splits \\
I10 & Rescue reversibility      & U   & Counterfactual activations exist for every example; no recovery is read \\
I11 & Onset coupling            & U   & Two states, before training and after; no sequence to locate an onset in \\
I12 & Offset coupling           & PC  & Behavior and structure both absent untrained; nothing in between \\
\midrule
\multicolumn{4}{@{}l}{\textit{External Validity}} \\
E1  & Intervention reach        & U   & One instrument throughout, so every agreement is the method with itself \\
E2  & Prompt generalization     & C   & Three prompt sets built differently, with the differences reported \\
E3  & Cross-task generalization & U   & One two-candidate setup throughout; generality claimed, not shown \\
E4  & Cross-model generalization& PC  & Transfers within the autoregressive class and fails outside it \\
E5  & Graded response           & PC  & Dose is the number of mediators, not the strength of the intervention \\
E6  & Novel prediction          & PC  & One prediction confirmed outside the model, one sharper one falsified \\
\midrule
\multicolumn{4}{@{}l}{\textit{Interpretive Validity}} \\
V1  & Level declaration         & C   & Structural-behavioral declared twice, and no claim sits outside it \\
V2  & Level-evidence match      & C   & Every headline claim is a property of the measured effect distribution \\
V3  & Alternative level         & PC  & The architecture rival is defeated; the decisive one is never raised \\
V4  & Unlicensed labeling       & C   & Components carry no names at all, only an index and an estimand \\
V5  & Scope declaration         & C   & Declared where each restriction binds, not collected at the end \\
\midrule
\multicolumn{4}{@{}l}{\textbf{Total:} 11 Confirmed, 14 Partially confirmed, 9 Untested, 2 Not applicable} \\
\multicolumn{4}{@{}l}{\textbf{Verdict: Causally Suggestive}} \\
\bottomrule
\end{longtable}}

\clearpage
\begin{table}[H]
\centering
\footnotesize
\caption{\textbf{Gender bias circuits.} Readings of the claim, from the version the authors state to the version the
field cites.}
\label{tab:gender-readings}
\renewcommand{\arraystretch}{1.18}
\begin{tabular}{@{}>{\raggedright\arraybackslash}p{3.3cm}l>{\raggedright\arraybackslash}p{4.4cm}@{}}
\toprule
\textbf{Reading} & \textbf{Verdict} & \textbf{Missing} \\
\midrule
\textbf{Primary.} Ten of 144 attention heads carry the indirect effect of a gender edit on the pronoun probability ratio, and neurons combine synergistically to produce it
  & Causally Suggestive
  & Specificity, which has no second outcome to register on (I4); necessity, which the design's substitution estimands cannot express (I1); correction for a greedy maximum over 144 candidates (M7); and dispersion on the head-level and neuron-level effects that carry the sparsity claim (M1) \\
\addlinespace[2pt]
These components mediate gender bias rather than gender processing
  & Underdetermined
  & A localized gender-competence mechanism to compute an overlap against; none is measured, and the definitionally gendered professions that would carry the test are excluded from the total-effect calculation (C4, V3, I5) \\
\addlinespace[2pt]
Training is what installs the sparse mediation structure
  & Underdetermined
  & A trajectory. The coupling is a single endpoint contrast against one untrained model whose total effects are small rather than absent --- 0.07 on WinoBias against GPT2-small's 0.25 (I12) --- with no intermediate checkpoint anywhere (I11) \\
\addlinespace[2pt]
The effect grows with model size and the pattern holds across model families
  & Disconfirmed
  & Nothing --- tested and failed: the neuron-level pattern does not transfer to the masked language models and the authors state they have no theory for the difference (E4, M6), and the size trend holds on the Winograd sets but not on Professions, where model size is not well correlated with total effect (E4) \\
\bottomrule
\end{tabular}
\end{table}

\begin{table}[H]
\centering
\footnotesize
\caption{\textbf{Gender bias circuits.} Verdict: \textbf{Causally Suggestive}. A dash means the tier is reached; a criterion at
Inconclusive, Untested or Disconfirmed blocks promotion.}
\label{tab:gender-verdict}
\renewcommand{\arraystretch}{1.18}
\begin{tabular}{@{}l>{\raggedright\arraybackslash}p{3.6cm}>{\raggedright\arraybackslash}p{4.4cm}@{}}
\toprule
\textbf{Tier} & \textbf{Requires} & \textbf{Missing} \\
\midrule
Proposed & C1--C2; one admissible measurement & --- \\
\addlinespace[2pt]
Causally Suggestive & Necessity (I1); baselines (M2) & \textbf{I1} not applicable: both estimands are substitutions, so no removal experiment exists in the design \\
\addlinespace[2pt]
Mechanistically Supported & Sufficiency (I2); reach (E1); specificity (I4) & \textbf{I4} unattempted and structurally so: the outcome is a two-element set, so nothing off-target can register; \textbf{E1} one intervention family throughout \\
\addlinespace[2pt]
Triangulated & Convergence (C3); replication (E2/E4); dissociation (I6) & \textbf{I6} unattempted; \textbf{E4} the neuron-level pattern fails outside the autoregressive families; \textbf{C3}: every number is an NIE from the same interchange intervention \\
\addlinespace[2pt]
Validated & Calibration (M1--M7); interpretive audit (V1--V5); reach (E2, E4) & \textbf{M7} unattempted: the head set is a greedy maximum over 144 candidates reported at its selected size; \textbf{M1} no dispersion on the effects carrying sparsity; \textbf{I8}, \textbf{I10}, \textbf{I11}, \textbf{E3} unattempted \\
\bottomrule
\end{tabular}
\end{table}


\claimheading{IOI Circuit}

\textbf{Source.} \emph{Interpretability in the Wild} \citep{wang2023interpretability}.

\textbf{Description.} Path patching over GPT-2 small identifies 26 attention heads in 7 classes---Previous Token, Duplicate Token, Induction, S-Inhibition, Name Mover, Negative Name Mover, Backup Name Mover---for sentences of the form ``When Mary and John went to the store, John gave a drink to \underline{\hspace{0.9cm}}''. Name Mover roles are supported at the weight level (OV copy score $>$95\%), the rest at the activation level. Component-level redundancy (Backup Name Movers) is recorded under minimality (I3) rather than necessity: a circuit whose parts compensate for one another is still necessary as a circuit. Description mode: implementational-functional.

\vspace{4pt}
{\footnotesize
\renewcommand{\arraystretch}{1.0}
\begin{longtable}{@{}llll>{\raggedright\arraybackslash}p{4.6cm}@{}}
\caption{\textbf{IOI circuit.} Full 36-criterion audit. Status: \textbf{C} = Confirmed,
\textbf{PC} = Partially confirmed, \textbf{U} = Untested, \textbf{I} = Inconclusive,
\textbf{D} = Disconfirmed, \textbf{N/A} = Not applicable.}\\
\toprule
\textbf{ID} & \textbf{Criterion} & \textbf{Status} & \textbf{Evidence} \\
\endfirsthead
\toprule
\textbf{ID} & \textbf{Criterion} & \textbf{Status} & \textbf{Evidence} \\
\endhead
\midrule
\multicolumn{4}{@{}l}{\textit{Construct Validity}} \\
C1  & Falsifiability            & C   & 26 heads named with roles; criteria stated as failable quantities \\
C2  & Structural plausibility   & PC  & Weight-level for two of seven classes; three omissions the authors name \\
C3  & Convergent validity       & PC  & Five evidence types, but only three share the mean-ablation primitive \\
C4  & Discriminant validity     & I   & 78\% head overlap with Colored Objects; opposite result elsewhere \\
C5  & Nomological validity      & PC  & Induction-head theory applied operationally; no formal derivation \\
C6  & Complementation validity  & U   & Seven named classes, knocked out one at a time; no trans comparison \\
\midrule
\multicolumn{4}{@{}l}{\textit{Measurement Validity}} \\
M1  & Reliability               & I   & SDs plotted for attention probabilities; none on faithfulness \\
M2  & Baseline separation       & PC  & Na\"ive circuit scores 0.1 against the full circuit's 0.46 \\
M3  & Stability                 & D   & Six method choices move the same quantity below 0\% to over 100\% \\
M4  & Calibration               & PC  & Logit difference has sign, null and baseline; faithfulness does not \\
M5  & Sensitivity               & U   & Unattempted; no planted circuit of known extent \\
M6  & Invariance                & D   & Faithfulness differs across the origin's own templates \\
M7  & Selection correction      & U   & N is unstated for the post-knockout sweep, the selection step itself \\
\midrule
\multicolumn{4}{@{}l}{\textit{Internal Validity}} \\
I1  & Necessity                 & PC  & Name Mover knockout costs only 5\%; the paper says so itself \\
I2  & Sufficiency               & PC  & Mean-ablating outside the circuit leaves 87\% of the logit difference \\
I3  & Minimality                & I   & Every node clears 1\%; Backup heads make components individually replaceable \\
I4  & Specificity               & I   & Verdict depends on whether head overlap or task effect is measured \\
I5  & Rival mechanism exclusion & I   & Greedy search finds knockout sets carrying 87\% of the behavior \\
I6  & Double dissociation       & U   & Unattempted; no published test located \\
I7  & Confound control          & PC  & Sequence length controlled twice; name frequency is not \\
I8  & Confounding sensitivity   & U   & Unattempted; no E-value or sensitivity analysis \\
I9  & Epistatic interaction     & PC  & Backup Name Movers are a discovered non-additivity \\
I10 & Rescue reversibility      & PC  & Direction reversed at origin; corrupt-then-restore run later \\
I11 & Onset coupling            & N/A & Untestable in GPT-2 small, a released checkpoint with no training run \\
I12 & Offset coupling           & PC  & A Pythia name-mover head loses the behavior late in training \\
\midrule
\multicolumn{4}{@{}l}{\textit{External Validity}} \\
E1  & Intervention reach        & I   & One ablation value at origin, three granularities; methods disagree \\
E2  & Prompt generalization     & D   & At $10^6$ clean/corrupted pairs, model and circuit diverge \\
E3  & Cross-task generalization & PC  & Repeated-random-token generalization at origin, one direction \\
E4  & Cross-model generalization& PC  & GPT-2 medium preliminary; not all heads attend to IO and S \\
E5  & Graded response           & PC  & Graded in circuit extent, never in intervention strength \\
E6  & Novel prediction          & C   & Duplicated-name prediction, with a matched control \\
\midrule
\multicolumn{4}{@{}l}{\textit{Interpretive Validity}} \\
V1  & Level declaration         & PC  & Circuits formally defined in \S2.1; description level not declared \\
V2  & Level-evidence match      & PC  & Topography and information flow supported; role semantics exceed it \\
V3  & Alternative level         & I   & The na\"ive circuit is simpler and not clearly worse \\
V4  & Unlicensed labeling       & PC  & ``Hedge'' for the Negative Name Movers, marked by the authors as speculation \\
V5  & Scope declaration         & C   & Limits declared in the abstract, \S1, \S3 and \S5 \\
\midrule
\multicolumn{4}{@{}l}{\textbf{Total:} 3 Confirmed, 17 Partially confirmed, 7 Inconclusive, 5 Untested, 3 Disconfirmed, 1 Not applicable} \\
\multicolumn{4}{@{}l}{\textbf{Verdict: Causally Suggestive}} \\
\bottomrule
\end{longtable}}

\clearpage
\begin{table}[H]
\centering
\footnotesize
\caption{\textbf{IOI circuit.} Readings of the claim, from the version the authors state to the version the
field cites.}
\label{tab:ioi-readings}
\renewcommand{\arraystretch}{1.18}
\begin{tabular}{@{}>{\raggedright\arraybackslash}p{3.3cm}l>{\raggedright\arraybackslash}p{4.4cm}@{}}
\toprule
\textbf{Reading} & \textbf{Verdict} & \textbf{Missing} \\
\midrule
\textbf{Primary.} A circuit comprising these 26 heads performs IOI in GPT-2 small on the origin templates
  & Causally Suggestive
  & Independent evidence (C3), separation from other circuits (I6), and robustness to method and prompt (M3, M6) \\
\addlinespace[2pt]
These 26 heads are \emph{the} mechanism, to the exclusion of others
  & Underdetermined
  & A test that discriminates among rival head sets; na\"ive and greedy-search circuits reach comparable faithfulness (I5) \\
\addlinespace[2pt]
The mechanism generalizes across prompts and models
  & Disconfirmed
  & Nothing --- tested and failed: model and circuit diverge at $10^6$ prompt pairs (E2), and faithfulness varies across the origin's own templates (M6) \\
\addlinespace[2pt]
The seven classes compose in the stated order, each performing the role its name asserts
  & Underdetermined
  & Weight-level accounts for four of seven classes (C2); role semantics currently exceed the evidence (V2) \\
\bottomrule
\end{tabular}
\end{table}

\begin{table}[H]
\centering
\footnotesize
\caption{\textbf{IOI circuit.} Verdict: \textbf{Causally Suggestive}. A dash means the tier is reached; a criterion at
Inconclusive, Untested or Disconfirmed blocks promotion.}
\label{tab:ioi-verdict}
\renewcommand{\arraystretch}{1.18}
\begin{tabular}{@{}l>{\raggedright\arraybackslash}p{3.6cm}>{\raggedright\arraybackslash}p{4.4cm}@{}}
\toprule
\textbf{Tier} & \textbf{Requires} & \textbf{Missing} \\
\midrule
Proposed & C1--C2; one admissible measurement & --- \\
\addlinespace[2pt]
Causally Suggestive & Necessity (I1); baselines (M2) & --- \\
\addlinespace[2pt]
Mechanistically Supported & Sufficiency (I2); reach (E1); specificity (I4) & \textbf{E1}, \textbf{I4} inconclusive \\
\addlinespace[2pt]
Triangulated & Convergence (C3); replication (E2/E4); dissociation (I6) & \textbf{C3}: the five families all descend from ablation; \textbf{E2} Disconfirmed; \textbf{I6} unattempted \\
\addlinespace[2pt]
Validated & Calibration (M1--M7); interpretive audit (V1--V5); reach (E2, E4) & \textbf{M3}, \textbf{M6}, \textbf{E2} Disconfirmed \\
\bottomrule
\end{tabular}
\end{table}


\claimheading{Othello World Model}

\textbf{Source.} \emph{Emergent World Representations: Exploring a Sequence Model Trained on a Synthetic Task} \citep{li2023othello}.

\textbf{Description.} A GPT variant trained from random initialization on Othello transcripts carries the board state in its activations. Nonlinear probes decode all 64 tiles where linear probes never dip below 20\% error, and 2,000 interventions move the top-$N$ predictions onto the counterfactual legal-move set (errors 0.12 and 0.06 against null baselines of 2.68 and 2.59), half of them on boards unreachable by legal play. Evaluated on: one 8-layer architecture trained twice, on championship games and on uniformly sampled legal moves, with both runs probed across all eight layers. Description mode: representational.

\vspace{4pt}
{\footnotesize
\renewcommand{\arraystretch}{1.0}
\begin{longtable}{@{}llll>{\raggedright\arraybackslash}p{4.6cm}@{}}
\caption{\textbf{Othello world model.} Full 36-criterion audit. Status: \textbf{C} = Confirmed,
\textbf{PC} = Partially confirmed, \textbf{U} = Untested, \textbf{I} = Inconclusive,
\textbf{D} = Disconfirmed, \textbf{N/A} = Not applicable.}\\
\toprule
\textbf{ID} & \textbf{Criterion} & \textbf{Status} & \textbf{Evidence} \\
\endfirsthead
\toprule
\textbf{ID} & \textbf{Criterion} & \textbf{Status} & \textbf{Evidence} \\
\endhead
\midrule
\multicolumn{4}{@{}l}{\textit{Construct Validity}} \\
C1  & Falsifiability            & C   & Intervention fixes the failing outcome before it is measured \\
C2  & Structural plausibility   & PC  & Layerwise argument and probe geometry; no weight-level account \\
C3  & Convergent validity       & C   & Probe, intervention and three metrics agree; two share one instrument \\
C4  & Discriminant validity     & PC  & Separated from memorisation and from a random net, not from a summary \\
C5  & Nomological validity      & PC  & Named rival theory and a literature; no theory of what a world model is \\
C6  & Complementation validity  & U   & Sixty-four tiles, intervened on one at a time throughout \\
\midrule
\multicolumn{4}{@{}l}{\textit{Measurement Validity}} \\
M1  & Reliability               & C   & Probe accuracies re-run 100 times; deviations in Tables 4 and 5 \\
M2  & Baseline separation       & C   & Untrained net, constant guess and null intervention floors all reported \\
M3  & Stability                 & PC  & Optimizer, ${\beta}$ and probe capacity swept; the probe target never varied \\
M4  & Calibration               & PC  & Intervention error has a null floor; probe accuracy has no task scale \\
M5  & Sensitivity               & PC  & A known-negative control only; no planted representation is recovered \\
M6  & Invariance                & C   & All eight layers, both datasets, both benchmarks, three metrics \\
M7  & Selection correction      & U   & Headline numbers are the maximum over swept Ls and ${\beta}$, uncorrected \\
\midrule
\multicolumn{4}{@{}l}{\textit{Internal Validity}} \\
I1  & Necessity                 & PC  & Substitution to a counterfactual state; nothing is removed \\
I2  & Sufficiency               & PC  & Editing the representation alone moves predictions, including off-distribution \\
I3  & Minimality                & N/A & One activation vector holds all 64 tiles; no components to prune \\
I4  & Specificity               & PC  & Off-target flips raised, mitigation swept, effects never counted \\
I5  & Rival mechanism exclusion & PC  & Memorisation and off-distribution rivals excluded; basis choice was not \\
I6  & Double dissociation       & U   & One representation, one behavior; no converse leg is available \\
I7  & Confound control          & PC  & Memorisation, distribution and probe capacity controlled; target is not \\
I8  & Confounding sensitivity   & U   & No sensitivity analysis, and no quantity that could serve as a bound \\
I9  & Epistatic interaction     & U   & Single-tile edits only; rule composition named as future work \\
I10 & Rescue reversibility      & PC  & Restoration observed as an obstacle, not designed as a rescue test \\
I11 & Onset coupling            & U   & Training-step axis available and explicitly declined \\
I12 & Offset coupling           & U   & Two models differ in competence; the coupling is never analyzed \\
\midrule
\multicolumn{4}{@{}l}{\textit{External Validity}} \\
E1  & Intervention reach        & U   & One intervention operator, reported under three outcome metrics \\
E2  & Prompt generalization     & C   & 1000 natural and 1000 unnatural cases; both training distributions \\
E3  & Cross-task generalization & PC  & Othello only; other games and natural language listed as future work \\
E4  & Cross-model generalization& C   & One architecture, two training sets; nothing cross-model at origin \\
E5  & Graded response           & PC  & Two knobs swept with degradation outside the optimum; no dose-response \\
E6  & Novel prediction          & C   & Unnatural boards predicted to work and did; saliency maps split two models \\
\midrule
\multicolumn{4}{@{}l}{\textit{Interpretive Validity}} \\
V1  & Level declaration         & PC  & Level named by the phrase used, and the phrase changes between sections \\
V2  & Level-evidence match      & PC  & Evidence matches decodable plus causal; the definition asks for more \\
V3  & Alternative level         & PC  & Memorisation alternative posed and refuted; state summary never posed \\
V4  & Unlicensed labeling       & PC  & ``World model'' is defined, and only a state summary is ever measured \\
V5  & Scope declaration         & C   & Synthetic scope declared in the title, §1.1 and the conclusion \\
\midrule
\multicolumn{4}{@{}l}{\textbf{Total:} 9 Confirmed, 18 Partially confirmed, 8 Untested, 1 Not applicable} \\
\multicolumn{4}{@{}l}{\textbf{Verdict: Causally Suggestive}} \\
\bottomrule
\end{longtable}}

\clearpage
\begin{table}[H]
\centering
\footnotesize
\caption{\textbf{Othello world model.} Readings of the claim, from the version the authors state to the version the
field cites.}
\label{tab:othello-readings}
\renewcommand{\arraystretch}{1.18}
\begin{tabular}{@{}>{\raggedright\arraybackslash}p{3.3cm}l>{\raggedright\arraybackslash}p{4.4cm}@{}}
\toprule
\textbf{Reading} & \textbf{Verdict} & \textbf{Missing} \\
\midrule
\textbf{Primary.} Othello-GPT carries a decodable representation of board state that its move predictions causally use
  & Causally Suggestive
  & Necessity by removal (I1), a second intervention operator (E1), a dissociation (I6), and a known-positive control, which the origin does not contain (M5) \\
\addlinespace[2pt]
The representation is a model of the process producing the sequences, which is the paper's own definition of a world model
  & Underdetermined
  & An experiment separating that from a decodable, causally-used state summary, which every result in the paper also satisfies (V4, V3); the rule composition that would distinguish them is measured once and fails on the OR across lines (V2, I9) \\
\addlinespace[2pt]
The board state is encoded nonlinearly
  & Underdetermined
  & A perturbation of the probe target, held at black/white/empty throughout and the one configuration the stability sweep never varied (M3, I5) \\
\addlinespace[2pt]
The account travels beyond Othello and beyond this architecture
  & Insufficient
  & The origin runs one task on one architecture trained twice, and lists other games and natural language as future work; E3 and E4 are scored on post-origin evidence \\
\bottomrule
\end{tabular}
\end{table}

\begin{table}[H]
\centering
\footnotesize
\caption{\textbf{Othello world model.} Verdict: \textbf{Causally Suggestive}. A dash means the tier is reached; a criterion at
Inconclusive, Untested or Disconfirmed blocks promotion.}
\label{tab:othello-verdict}
\renewcommand{\arraystretch}{1.18}
\begin{tabular}{@{}l>{\raggedright\arraybackslash}p{3.6cm}>{\raggedright\arraybackslash}p{4.4cm}@{}}
\toprule
\textbf{Tier} & \textbf{Requires} & \textbf{Missing} \\
\midrule
Proposed & C1--C2; one admissible measurement & --- \\
\addlinespace[2pt]
Causally Suggestive & Necessity (I1); baselines (M2) & --- \\
\addlinespace[2pt]
Mechanistically Supported & Sufficiency (I2); reach (E1); specificity (I4) & \textbf{E1} unattempted: one intervention operator throughout, reported under three agreeing outcome metrics; \textbf{I4} off-target tile flips acknowledged and never counted \\
\addlinespace[2pt]
Triangulated & Convergence (C3); replication (E2/E4); dissociation (I6) & \textbf{I6} unattempted, one representation and one behavior; \textbf{C3}: probe and intervention are one instrument used twice; \textbf{E4} rests entirely on post-origin work \\
\addlinespace[2pt]
Validated & Calibration (M1--M7); interpretive audit (V1--V5); reach (E2, E4) & \textbf{M7} unattempted, and both headline errors are the best cell of a swept $L_s \times \beta$ grid; \textbf{M5} supplies a known-negative only; \textbf{I8}, \textbf{I9}, \textbf{I11}, \textbf{I12} unattempted \\
\bottomrule
\end{tabular}
\end{table}


\claimheading{SAE Features as Units of Analysis}

\textbf{Source.} \emph{Sparse Autoencoders Find Highly Interpretable Features} \citep{cunningham2024sparse}.

\textbf{Description.} Sparse autoencoders are trained on residual-stream activations to recover an overcomplete dictionary of directions, evaluated two ways: automated interpretability scoring of 150 features per method against six baselines (default basis, random directions, PCA, ICA, top-$K$ PCA, top-$K$ ICA), and activation patching on 50 IOI data points. Evaluated on: Pythia-70M and 410M at origin, GPT-2 small and Gemma 2 2B in the follow-up. Being a method-level claim, these are the systems the evidence was gathered on rather than one audited model. Description mode: representational.

\vspace{4pt}
{\footnotesize
\renewcommand{\arraystretch}{1.0}
\begin{longtable}{@{}llll>{\raggedright\arraybackslash}p{4.6cm}@{}}
\caption{\textbf{SAE features.} Full 36-criterion audit. Status: \textbf{C} = Confirmed,
\textbf{PC} = Partially confirmed, \textbf{U} = Untested, \textbf{I} = Inconclusive,
\textbf{D} = Disconfirmed, \textbf{N/A} = Not applicable.}\\
\toprule
\textbf{ID} & \textbf{Criterion} & \textbf{Status} & \textbf{Evidence} \\
\endfirsthead
\toprule
\textbf{ID} & \textbf{Criterion} & \textbf{Status} & \textbf{Evidence} \\
\endhead
\midrule
\multicolumn{4}{@{}l}{\textit{Construct Validity}} \\
C1  & Falsifiability            & C   & Directional prediction against four named baselines; partly fails in layer 4 \\
C2  & Structural plausibility   & PC  & Predicted by superposition; MLP training fails, attention not attempted \\
C3  & Convergent validity       & PC  & Two evidence types; the concurrent-origin pairing comes from the follow-up, not here \\
C4  & Discriminant validity     & PC  & Scored against the origin's operational construct, not the field's \\
C5  & Nomological validity      & PC  & Fits superposition; the theory does not underwrite the method \\
C6  & Complementation validity  & U   & Feature splitting is the open question; no test of atomicity is run \\
\midrule
\multicolumn{4}{@{}l}{\textit{Measurement Validity}} \\
M1  & Reliability               & I   & One training run per configuration; no repetition reported \\
M2  & Baseline separation       & PC  & SAE beats four baselines; the two negative controls do not separate from each other \\
M3  & Stability                 & I   & Sparsity sweep gives a smooth tradeoff, not a stable optimum \\
M4  & Calibration               & PC  & Interpretability score is a correlation with a meaningful zero \\
M5  & Sensitivity               & U   & Unattempted; no planted ground-truth features \\
M6  & Invariance                & PC  & All layers of Pythia-70M, five expansion ratios, two scoring regimes \\
M7  & Selection correction      & U   & A feature scoring fewer than 20 varying fragments is skipped; the count is unreported \\
\midrule
\multicolumn{4}{@{}l}{\textit{Internal Validity}} \\
I1  & Necessity                 & PC  & One feature ablated for behavior; a full previous-layer sweep reads out to features \\
I2  & Sufficiency               & PC  & Interchange intervention on IOI; counterfactual features written in \\
I3  & Minimality                & PC  & ACDC orders features; the next feature usually adds little \\
I4  & Specificity               & I   & Single-feature ablation moves 12,000 logits; left unanalyzed \\
I5  & Rival mechanism exclusion & PC  & Rival methods excluded: PCA, ICA, top-$K$, non-sparse dictionary \\
I6  & Double dissociation       & U   & Unattempted; no second decomposition shown intact \\
I7  & Confound control          & PC  & Three designed controls, including an $\alpha=0$ dictionary \\
I8  & Confounding sensitivity   & U   & Unattempted; no E-value or sensitivity analysis \\
I9  & Epistatic interaction     & PC  & Assumed away explicitly; ACDC treats features as a flat graph \\
I10 & Rescue reversibility      & U   & Unattempted, though patching is reversible by construction \\
I11 & Onset coupling            & U   & Unattempted, and cheap: Pythia ships pretraining checkpoints \\
I12 & Offset coupling           & U   & Unattempted; the nearest result is at the extreme rather than along a trajectory \\
\midrule
\multicolumn{4}{@{}l}{\textit{External Validity}} \\
E1  & Intervention reach        & PC  & Two forms: activation patching and less-than-rank-one ablation \\
E2  & Prompt generalization     & PC  & Cross-corpus by construction: trained on the Pile, scored on OpenWebText \\
E3  & Cross-task generalization & U   & Unattempted; generalization stated as an expectation \\
E4  & Cross-model generalization& C   & Transfers everywhere tried: Pythia, GPT-2 small, Gemma 2 2B \\
E5  & Graded response           & PC  & Two graded curves: KL against features patched, edit size against KL \\
E6  & Novel prediction          & PC  & The IOI experiment is a prediction in the weak sense \\
\midrule
\multicolumn{4}{@{}l}{\textit{Interpretive Validity}} \\
V1  & Level declaration         & PC  & Formal object declared precisely; ``feature'' carries more than it \\
V2  & Level-evidence match      & PC  & Evidence supports the decomposition claim the origin makes \\
V3  & Alternative level         & PC  & Raised at origin and not pursued: no single correct decomposition \\
V4  & Unlicensed labeling       & PC  & ``Monosemantic'' projects semantics; the origin hedges consistently \\
V5  & Scope declaration         & C   & Both papers declare limits and quantify them \\
\midrule
\multicolumn{4}{@{}l}{\textbf{Total:} 3 Confirmed, 21 Partially confirmed, 3 Inconclusive, 9 Untested} \\
\multicolumn{4}{@{}l}{\textbf{Verdict: Causally Suggestive}} \\
\bottomrule
\end{longtable}}

\clearpage
\begin{table}[H]
\centering
\footnotesize
\caption{\textbf{SAE features.} Readings of the claim, from the version the authors state to the version the
field cites.}
\label{tab:sae-readings}
\renewcommand{\arraystretch}{1.18}
\begin{tabular}{@{}>{\raggedright\arraybackslash}p{3.3cm}l>{\raggedright\arraybackslash}p{4.4cm}@{}}
\toprule
\textbf{Reading} & \textbf{Verdict} & \textbf{Missing} \\
\midrule
\textbf{Primary.} Sparse dictionary learning recovers directions more interpretable than PCA, ICA and the neuron basis, and causally usable
  & Causally Suggestive
  & Specificity (I4), tested and inconclusive; then independent convergence (C3), separation from other decompositions (I6), and calibration of the scoring instrument (M1, M3, M5, M7) \\
\addlinespace[2pt]
SAE features are atomic, canonical units of the model
  & Disconfirmed
  & Nothing --- tested and failed: meta-SAEs decompose latents further (C4), and the instrument cannot distinguish a trained transformer from a random one \citep{heap2025sparse} \\
\addlinespace[2pt]
SAE features localize causation better than neurons
  & Disconfirmed
  & Nothing --- tested and failed: \citet{mueller2025mib} report no causal-localization advantage over the neuron basis \\
\bottomrule
\end{tabular}
\end{table}

\begin{table}[H]
\centering
\footnotesize
\caption{\textbf{SAE features.} Verdict: \textbf{Causally Suggestive}. A dash means the tier is reached; a criterion at
Inconclusive, Untested or Disconfirmed blocks promotion.}
\label{tab:sae-verdict}
\renewcommand{\arraystretch}{1.18}
\begin{tabular}{@{}l>{\raggedright\arraybackslash}p{3.6cm}>{\raggedright\arraybackslash}p{4.4cm}@{}}
\toprule
\textbf{Tier} & \textbf{Requires} & \textbf{Missing} \\
\midrule
Proposed & C1--C2; one admissible measurement & --- \\
\addlinespace[2pt]
Causally Suggestive & Necessity (I1); baselines (M2) & --- \\
\addlinespace[2pt]
Mechanistically Supported & Sufficiency (I2); reach (E1); specificity (I4) & \textbf{E1}, \textbf{I4} inconclusive \\
\addlinespace[2pt]
Triangulated & Convergence (C3); replication (E2/E4); dissociation (I6) & \textbf{C3}: later work reaches opposite conclusions on utility; \textbf{I6} unattempted \\
\addlinespace[2pt]
Validated & Calibration (M1--M7); interpretive audit (V1--V5); reach (E2, E4) & \textbf{M1}, \textbf{M3} Inconclusive; \textbf{M5}, \textbf{M7} unattempted \\
\bottomrule
\end{tabular}
\end{table}


\claimheading{Probing Classifiers as Mechanistic Evidence}

\textbf{Source.} \emph{Probing Classifiers: Promises, Shortcomings, and Advances} \citep{belinkov2022probing}.

\textbf{Description.} A classifier $g$ is trained on a model's intermediate output $f_l(x)$ to predict a property $z$, and its accuracy is reported as evidence that the model has learned information relevant to $z$. The audited object is the method rather than one experiment: the anchor is a review that runs no experiments of its own, so the evidence sits in the studies it formalizes. Evaluated on: every architecture family the method has been applied to, from static word embeddings to transformers. Description mode: representational.

\vspace{4pt}
{\footnotesize
\renewcommand{\arraystretch}{1.0}
\begin{longtable}{@{}llll>{\raggedright\arraybackslash}p{4.6cm}@{}}
\caption{\textbf{Probing classifiers.} Full 36-criterion audit. Status: \textbf{C} = Confirmed,
\textbf{PC} = Partially confirmed, \textbf{U} = Untested, \textbf{I} = Inconclusive,
\textbf{D} = Disconfirmed, \textbf{N/A} = Not applicable.}\\
\toprule
\textbf{ID} & \textbf{Criterion} & \textbf{Status} & \textbf{Evidence} \\
\endfirsthead
\toprule
\textbf{ID} & \textbf{Criterion} & \textbf{Status} & \textbf{Evidence} \\
\endhead
\midrule
\multicolumn{4}{@{}l}{\textit{Construct Validity}} \\
C1  & Falsifiability            & C   & The prediction can fail and did, in work the anchor reports side by side \\
C2  & Structural plausibility   & PC  & Every symbol a probing result depends on is named and typed \\
C3  & Convergent validity       & PC  & Probe families agree under accuracy and reverse under other measures \\
C4  & Discriminant validity     & D   & Two neighbors the instrument must separate: probe memorisation, random features \\
C5  & Nomological validity      & PC  & Tied to mutual information, description length and information gain \\
C6  & Complementation validity  & N/A & The audited object is a method, not a construct with named parts \\
\midrule
\multicolumn{4}{@{}l}{\textit{Measurement Validity}} \\
M1  & Reliability               & U   & A review reports no variance; the anchor states what it requires of studies \\
M2  & Baseline separation       & C   & The criterion the literature solved; three control families formalized \\
M3  & Stability                 & I   & Ordering survives under accuracy and inverts under an alternative measure \\
M4  & Calibration               & D   & The shortcomings section opens by denying the measurement has a scale \\
M5  & Sensitivity               & PC  & Skylines and floors exist; no planted property is ever recovered \\
M6  & Invariance                & PC  & Invariance is what probing measures, so the instrument is the question \\
M7  & Selection correction      & U   & Three control families catalogued; the omission from the catalogue is the gap \\
\midrule
\multicolumn{4}{@{}l}{\textit{Internal Validity}} \\
I1  & Necessity                 & I   & Standard probing performs no removal; removal studies disagree four ways \\
I2  & Sufficiency               & N/A & Probing is a map out of the representation, with no path writing back in \\
I3  & Minimality                & U   & The unit is a whole intermediate output; no operation removes part of one \\
I4  & Specificity               & D   & A control dataset holds the property constant and probing fails on it \\
I5  & Rival mechanism exclusion & PC  & Control tasks are solvable only by memorisation, and separate cleanly \\
I6  & Double dissociation       & U   & Needs two properties and two behaviors; the framework supplies one of each \\
I7  & Confound control          & PC  & Probe capacity, information in a random baseline and task difficulty \\
I8  & Confounding sensitivity   & U   & The anchor answers confounding with designed controls, never with a bound \\
I9  & Epistatic interaction     & U   & Non-additivity needs two manipulable components; the object is one output \\
I10 & Rescue reversibility      & U   & Every catalogued intervention runs one way, projecting or training out \\
I11 & Onset coupling            & U   & The original model is taken as given and trained once \\
I12 & Offset coupling           & U   & The converse needs the training sequence onset coupling also lacks \\
\midrule
\multicolumn{4}{@{}l}{\textit{External Validity}} \\
E1  & Intervention reach        & I   & Probing has no intervention of its own; those applied to it disagree \\
E2  & Prompt generalization     & PC  & One line of work reports the same property across several datasets \\
E3  & Cross-task generalization & PC  & Transfers to any property with an annotated dataset, which is the appeal \\
E4  & Cross-model generalization& C   & Portability is established without qualification, from static embeddings on \\
E5  & Graded response           & N/A & No intervention, so no strength to dial and no response to grade \\
E6  & Novel prediction          & PC  & Four design changes followed probing results and then worked \\
\midrule
\multicolumn{4}{@{}l}{\textit{Interpretive Validity}} \\
V1  & Level declaration         & PC  & The formalism fixes which objects a probing result is about \\
V2  & Level-evidence match      & D   & The evidence supports extractability; the claim made is representation \\
V3  & Alternative level         & C   & Every rival reading of a good score is given a measure of its own \\
V4  & Unlicensed labeling       & PC  & ``The model knows z'' projects an epistemic relation onto decodability \\
V5  & Scope declaration         & C   & A scope declaration end to end, conceding limits before any content \\
\midrule
\multicolumn{4}{@{}l}{\textbf{Total:} 5 Confirmed, 12 Partially confirmed, 3 Inconclusive, 9 Untested, 4 Disconfirmed, 3 Not applicable} \\
\multicolumn{4}{@{}l}{\textbf{Verdict: Proposed}} \\
\bottomrule
\end{longtable}}

\clearpage
\begin{table}[H]
\centering
\footnotesize
\caption{\textbf{Probing classifiers.} Readings of the claim, from the version the authors state to the version the
field cites.}
\label{tab:probing-readings}
\renewcommand{\arraystretch}{1.18}
\begin{tabular}{@{}>{\raggedright\arraybackslash}p{3.3cm}l>{\raggedright\arraybackslash}p{4.4cm}@{}}
\toprule
\textbf{Reading} & \textbf{Verdict} & \textbf{Missing} \\
\midrule
\textbf{Primary.} A probing result licenses that property $z$ is extractable from the representation by a classifier of the stated capacity
  & Proposed
  & Necessity: the method performs no removal of its own, and where removal was performed the studies the anchor reports disagree four ways (I1, E1) \\
\addlinespace[2pt]
A good probe score shows that the model represents $z$
  & Disconfirmed
  & Nothing --- tested and failed: control tasks attribute most of a nonlinear probe's accuracy to the probe itself, and even random features decode the property (C4) \\
\addlinespace[2pt]
A good probe score shows that the model uses $z$ in producing its output
  & Disconfirmed
  & Nothing --- tested and failed: a control dataset holds $z$ non-discriminative for the original task and the probe recovers it anyway, so decodability does not localize to the task (I4) \\
\addlinespace[2pt]
Interventions on a probe-identified direction show which features the model uses
  & Underdetermined
  & A result that decides the four-way disagreement among the intervention studies of \S4.3 (I1, E1). This is the broad scope, under which \textbf{I2} and \textbf{E5} are unattempted; under probing as \S2 defines it, a read-out with no path writing back in, they are not applicable \\
\bottomrule
\end{tabular}
\end{table}

\begin{table}[H]
\centering
\footnotesize
\caption{\textbf{Probing classifiers.} Verdict: \textbf{Proposed}. A dash means the tier is reached; a criterion at
Inconclusive, Untested or Disconfirmed blocks promotion.}
\label{tab:probing-verdict}
\renewcommand{\arraystretch}{1.18}
\begin{tabular}{@{}l>{\raggedright\arraybackslash}p{3.6cm}>{\raggedright\arraybackslash}p{4.4cm}@{}}
\toprule
\textbf{Tier} & \textbf{Requires} & \textbf{Missing} \\
\midrule
Proposed & C1--C2; one admissible measurement & --- \\
\addlinespace[2pt]
Causally Suggestive & Necessity (I1); baselines (M2) & \textbf{I1} inconclusive: the anchor sets Giulianelli's positive conclusion against Elazar's negative one and prefers neither \\
\addlinespace[2pt]
Mechanistically Supported & Sufficiency (I2); reach (E1); specificity (I4) & \textbf{I4} Disconfirmed; \textbf{E1} inconclusive; \textbf{I2} not applicable: a read-out has no write path \\
\addlinespace[2pt]
Triangulated & Convergence (C3); replication (E2/E4); dissociation (I6) & \textbf{I6} unattempted: one property and one behavior per experiment; \textbf{C3}: probe families agree under accuracy and reverse under selectivity \\
\addlinespace[2pt]
Validated & Calibration (M1--M7); interpretive audit (V1--V5); reach (E2, E4) & \textbf{M4}, \textbf{V2} Disconfirmed; \textbf{M3} inconclusive; \textbf{M1}, \textbf{M7} unattempted \\
\bottomrule
\end{tabular}
\end{table}


\claimheading{Induction Heads: General In-Context Learning}

\textbf{Source.} \emph{In-context Learning and Induction Heads} \citep{olsson2022context}.

\textbf{Description.} The two-head composition that performs in-context copying is advanced as the mechanistic source of in-context learning in general. In-context learning is operationalized as the difference in loss between the 500th and the 50th token of the context, and the claim rests on induction-head formation, the jump in that quantity, a loss bump and a PCA pivot all arriving together in a $2.5$--$5\times10^{9}$-token window. Evaluated on: the same 34 decoder-only models, with ablation confined to the twelve small models. Description mode: implementational-functional, carried to a computational-level capability by Argument 6.

\vspace{4pt}
{\footnotesize
\renewcommand{\arraystretch}{1.0}
\begin{longtable}{@{}llll>{\raggedright\arraybackslash}p{4.6cm}@{}}
\caption{\textbf{Induction heads: general in-context learning.} Full 36-criterion audit. Status: \textbf{C} = Confirmed,
\textbf{PC} = Partially confirmed, \textbf{U} = Untested, \textbf{I} = Inconclusive,
\textbf{D} = Disconfirmed, \textbf{N/A} = Not applicable.}\\
\toprule
\textbf{ID} & \textbf{Criterion} & \textbf{Status} & \textbf{Evidence} \\
\endfirsthead
\toprule
\textbf{ID} & \textbf{Criterion} & \textbf{Status} & \textbf{Evidence} \\
\endhead
\midrule
\multicolumn{4}{@{}l}{\textit{Construct Validity}} \\
C1  & Falsifiability            & C   & Argument 6 names the falsifier precisely, whatever the eventual outcome \\
C2  & Structural plausibility   & PC  & Weight-level for small attention-only models; the extension is analogy \\
C3  & Convergent validity       & PC  & Six lines of evidence, but at scale the paper's own table reads correlational \\
C4  & Discriminant validity     & D   & The title says general in-context learning; the measure is loss at token 500 \\
C5  & Nomological validity      & PC  & Connected in three directions, mostly by analogy rather than derivation \\
C6  & Complementation validity  & U   & The composition-head alternative is what a trans test would settle \\
\midrule
\multicolumn{4}{@{}l}{\textit{Measurement Validity}} \\
M1  & Reliability               & U   & One run per model, and at scale the evidence is timing across fifteen snapshots \\
M2  & Baseline separation       & PC  & The one-layer model is a clean negative for all four signature quantities \\
M3  & Stability                 & PC  & Free constants perturbed and the between-model picture survives \\
M4  & Calibration               & I   & Validated evaluators, but footnote 17 documents a sign-inverting artifact \\
M5  & Sensitivity               & U   & The only known-positive is literal copying in a two-layer model \\
M6  & Invariance                & PC  & The signature recurs across four series; a correlate's invariance is not the claim's \\
M7  & Selection correction      & U   & Top scorers examined on the data that ranked them, and the sample is biased \\
\midrule
\multicolumn{4}{@{}l}{\textit{Internal Validity}} \\
I1  & Necessity                 & I   & Ablation carries the metric in twelve models and the construct nowhere \\
I2  & Sufficiency               & U   & A semi-empirical argument stands in; nothing reconstructs the behavior \\
I3  & Minimality                & U   & Every ablation marginal and single-head, in the redundancy regime that defeats it \\
I4  & Specificity               & PC  & Ablating non-induction heads raises the score, which runs against the claim \\
I5  & Rival mechanism exclusion & D   & The composition rival is named twice by the authors and excluded neither time \\
I6  & Double dissociation       & U   & No second head population is defined that could be ablated against a second outcome \\
I7  & Confound control          & PC  & Exogenous confounds located outside the window; the shared latent is untouched \\
I8  & Confounding sensitivity   & U   & Every caution is verbal; the leading alternative is left unquantified \\
I9  & Epistatic interaction     & PC  & Composition is the subject and no joint ablation measures it \\
I10 & Rescue reversibility      & U   & The clean patterns are already cached; the restore costs one pass and is not run \\
I11 & Onset coupling            & C   & Four quantities move together in one token window, and an intervention shifts them \\
I12 & Offset coupling           & D   & The manipulation runs one way only; nothing suppresses formation and re-tests \\
\midrule
\multicolumn{4}{@{}l}{\textit{External Validity}} \\
E1  & Intervention reach        & PC  & Two levers of different kinds, neither reaching the models the claim is about \\
E2  & Prompt generalization     & C   & The defining stimulus is off-distribution by construction \\
E3  & Cross-task generalization & PC  & Three heads in one model on three tasks, labeled Plausibility by the authors \\
E4  & Cross-model generalization& D   & The signature holds across 34 models; nothing distinguishes the heads at scale \\
E5  & Graded response           & PC  & Continuous scores throughout and an all-or-nothing intervention \\
E6  & Novel prediction          & C   & The smeared-key architecture predicted the shift in advance, and it moved \\
\midrule
\multicolumn{4}{@{}l}{\textit{Interpretive Validity}} \\
V1  & Level declaration         & C   & Every sub-claim graded in a table, by reach and by model class \\
V2  & Level-evidence match      & D   & The assertion is mechanistic and unbounded where the tables say correlational \\
V3  & Alternative level         & PC  & The shared-latent alternative is stated precisely twice and tested never \\
V4  & Unlicensed labeling       & PC  & Both loaded terms hedged, and neither contrastively tested \\
V5  & Scope declaration         & C   & Scoped structurally, with the missing evidence graded rather than hedged \\
\midrule
\multicolumn{4}{@{}l}{\textbf{Total:} 6 Confirmed, 14 Partially confirmed, 2 Inconclusive, 9 Untested, 5 Disconfirmed} \\
\multicolumn{4}{@{}l}{\textbf{Verdict: Disconfirmed}} \\
\bottomrule
\end{longtable}}

\clearpage
\begin{table}[H]
\centering
\footnotesize
\caption{\textbf{Induction heads: general in-context learning.} Readings of the claim, from the version the authors state to the version the
field cites.}
\label{tab:induction_broad-readings}
\renewcommand{\arraystretch}{1.18}
\begin{tabular}{@{}>{\raggedright\arraybackslash}p{3.3cm}l>{\raggedright\arraybackslash}p{4.4cm}@{}}
\toprule
\textbf{Reading} & \textbf{Verdict} & \textbf{Missing} \\
\midrule
\textbf{Primary.} Induction heads are the mechanistic source of general in-context learning in transformers of any size
  & Disconfirmed
  & Nothing --- tested and failed: no ablation runs above the twelve small models, so nothing at scale separates induction heads from whatever else forms alongside them (E4, I5), and the adopted measure does not separate general in-context learning from the few-shot accuracy the field reads it as (C4) \\
\addlinespace[2pt]
The phase change at which induction heads form is when in-context learning arrives
  & Underdetermined
  & A test that separates the rival the authors state in Argument 1 --- that the phase change is when layer composition becomes available, forming induction heads and other composition mechanisms together --- from the causal reading (I5, I7, V3) \\
\addlinespace[2pt]
The share of in-context learning induction heads account for is the majority
  & Insufficient
  & No admissible measurement of the share: the paper's summary table asserts it while the semi-empirical argument offered for it states no fraction (I2), the ablation is single-head and all-or-nothing (E5), and the headline metric is written down five times with four of them reversing its sign (M4) \\
\addlinespace[2pt]
The copying the ablations cover is the literal case of a general retrieval mechanism
  & Underdetermined
  & The bridge is the in-context nearest-neighbor analogy, argued rather than derived (C2), and the paper's own strength table labels its large-model rows correlational and Argument 6 analogy (V2) \\
\bottomrule
\end{tabular}
\end{table}

\begin{table}[H]
\centering
\footnotesize
\caption{\textbf{Induction heads: general in-context learning.} Verdict: \textbf{Disconfirmed}. A dash means the tier is reached; a criterion at
Inconclusive, Untested or Disconfirmed blocks promotion.}
\label{tab:induction_broad-verdict}
\renewcommand{\arraystretch}{1.18}
\begin{tabular}{@{}l>{\raggedright\arraybackslash}p{3.6cm}>{\raggedright\arraybackslash}p{4.4cm}@{}}
\toprule
\textbf{Tier} & \textbf{Requires} & \textbf{Missing} \\
\midrule
Proposed & C1--C2; one admissible measurement & --- \\
\addlinespace[2pt]
Causally Suggestive & Necessity (I1); baselines (M2) & \textbf{I1} inconclusive: necessity holds for the metric in the twelve small models, is called only suggestive once MLPs are present, and is untested above them \\
\addlinespace[2pt]
Mechanistically Supported & Sufficiency (I2); reach (E1); specificity (I4) & \textbf{I2} unattempted: nothing reconstructs in-context learning from induction heads, and the argument offered instead states no fraction \\
\addlinespace[2pt]
Triangulated & Convergence (C3); replication (E2/E4); dissociation (I6) & \textbf{E4} Disconfirmed: the causal evidence never leaves the twelve small models; \textbf{C3}: only two of the six lines of evidence are interventional and both stop there \\
\addlinespace[2pt]
Validated & Calibration (M1--M7); interpretive audit (V1--V5); reach (E2, E4) & \textbf{V2}, \textbf{E4} Disconfirmed; \textbf{M4} inconclusive; \textbf{M1}, \textbf{M5}, \textbf{M7} unattempted, and the head pool is ranked and then sampled towards its high scorers \\
\bottomrule
\end{tabular}
\end{table}


\claimheading{Knowledge Neurons}

\textbf{Source.} \emph{Knowledge Neurons in Pretrained Transformers} \citep{dai2022knowledge}.

\textbf{Description.} Integrated gradients over the intermediate activations of a feed-forward layer attribute a relational fact to about four neurons, read as the value slots of the key--value memory that layer implements. Zeroing those activations lowers the correct-answer probability by 29.03\% and doubling them raises it by 31.17\%, against $-1.47\%$ and $-1.27\%$ for a count-matched baseline. Evaluated on: BERT-base-cased at one released checkpoint, on 253,448 ParaRel cloze prompts covering 27,738 facts. Description mode: implementational-functional.

\vspace{4pt}
{\footnotesize
\renewcommand{\arraystretch}{1.0}
\begin{longtable}{@{}llll>{\raggedright\arraybackslash}p{4.6cm}@{}}
\caption{\textbf{Knowledge neurons.} Full 36-criterion audit. Status: \textbf{C} = Confirmed,
\textbf{PC} = Partially confirmed, \textbf{U} = Untested, \textbf{I} = Inconclusive,
\textbf{D} = Disconfirmed, \textbf{N/A} = Not applicable.}\\
\toprule
\textbf{ID} & \textbf{Criterion} & \textbf{Status} & \textbf{Evidence} \\
\endfirsthead
\toprule
\textbf{ID} & \textbf{Criterion} & \textbf{Status} & \textbf{Evidence} \\
\endhead
\midrule
\multicolumn{4}{@{}l}{\textit{Construct Validity}} \\
C1  & Falsifiability            & C   & Signed predictions run both ways on 34 relations, against a matched control \\
C2  & Structural plausibility   & C   & Eq. (3) beside Eq. (2): the same key-value operation under another nonlinearity \\
C3  & Convergent validity       & PC  & One attributor against one contrast, agreeing on one coarse property \\
C4  & Discriminant validity     & D   & The competitor is named, removed by a filter, and the same machinery moves it \\
C5  & Nomological validity      & PC  & Two theories say where to look; neither is tested as a network \\
C6  & Complementation validity  & U   & Relation sets separate at identification; no joint suppression is run \\
\midrule
\multicolumn{4}{@{}l}{\textit{Measurement Validity}} \\
M1  & Reliability               & U   & Every quantity a mean reported once, with no interval or repeated run \\
M2  & Baseline separation       & C   & Activation magnitude through the identical pipeline, succeeding 0.0\% of the time \\
M3  & Stability                 & U   & Three free settings stated and none of them perturbed \\
M4  & Calibration               & PC  & Outcomes calibrated; the selection quantity is driven to a target range \\
M5  & Sensitivity               & U   & One known-negative control and no case where the answer is known in advance \\
M6  & Invariance                & PC  & Per-relation spread across 34 relations; the generalization claim has no experiment \\
M7  & Selection correction      & U   & Selection at four places, disclosed at all four and corrected at none \\
\midrule
\multicolumn{4}{@{}l}{\textit{Internal Validity}} \\
I1  & Necessity                 & PC  & 29.03\% against the control's 1.47\%, consistent across relations but a decrement \\
I2  & Sufficiency               & PC  & Doubling gives +31.17\% and rewriting succeeds 34.4\%; neither reaches sufficiency \\
I3  & Minimality                & U   & Set size is an input, held near four neurons before any effect is measured \\
I4  & Specificity               & I   & Holds at identification and erasure, fails on other knowledge: results point both ways \\
I5  & Rival mechanism exclusion & U   & The one alternative ruled out is about method; the FFN framing forecloses the rest \\
I6  & Double dissociation       & U   & One mechanism localized and one behavior measured, so neither arm exists \\
I7  & Confound control          & PC  & Wording controlled by requiring recurrence across nine templates; frequency is not \\
I8  & Confounding sensitivity   & U   & Unattempted; nothing bounds an unmeasured confounder \\
I9  & Epistatic interaction     & U   & Every manipulation hits the whole set at once, so no contribution separates \\
I10 & Rescue reversibility      & U   & The damage is analytic and invertible, and the restore is never run \\
I11 & Onset coupling            & N/A & One released checkpoint, so there is no sequence to locate an onset in \\
I12 & Offset coupling           & N/A & One checkpoint supplies no interval over which either could go \\
\midrule
\multicolumn{4}{@{}l}{\textit{External Validity}} \\
E1  & Intervention reach        & PC  & Three forms agreeing in direction, all downstream of one attributor \\
E2  & Prompt generalization     & C   & Survival across roughly nine templates is built into the identification \\
E3  & Cross-task generalization & U   & Single-word cloze throughout, named first among the authors' limitations \\
E4  & Cross-model generalization& PC  & One model, with generalization asserted and no experiment behind it \\
E5  & Graded response           & PC  & Sign reverses between zero and double; nothing reported between them \\
E6  & Novel prediction          & C   & Unseen text fires the neurons where it expresses the fact and not otherwise \\
\midrule
\multicolumn{4}{@{}l}{\textit{Interpretive Validity}} \\
V1  & Level declaration         & PC  & The unit is declared exactly and the level is not; storage and expression run together \\
V2  & Level-evidence match      & D   & The title claims storage; the measurement is correlation with expression \\
V3  & Alternative level         & I   & The method alternative is rejected; the storage-versus-expression reading is untouched \\
V4  & Unlicensed labeling       & PC  & The name arrives in the sentence that introduces the method, before the measurement \\
V5  & Scope declaration         & PC  & Four limits declared; the one asserted past is generalization beyond BERT \\
\midrule
\multicolumn{4}{@{}l}{\textbf{Total:} 5 Confirmed, 13 Partially confirmed, 2 Inconclusive, 12 Untested, 2 Disconfirmed, 2 Not applicable} \\
\multicolumn{4}{@{}l}{\textbf{Verdict: Disconfirmed}} \\
\bottomrule
\end{longtable}}

\clearpage
\begin{table}[H]
\centering
\footnotesize
\caption{\textbf{Knowledge neurons.} Readings of the claim, from the version the authors state to the version the
field cites.}
\label{tab:knowledge_neurons-readings}
\renewcommand{\arraystretch}{1.18}
\begin{tabular}{@{}>{\raggedright\arraybackslash}p{3.3cm}l>{\raggedright\arraybackslash}p{4.4cm}@{}}
\toprule
\textbf{Reading} & \textbf{Verdict} & \textbf{Missing} \\
\midrule
\textbf{Primary.} These roughly four feed-forward neurons store the relational fact
  & Disconfirmed
  & Nothing --- tested and failed: all three of the paper's own summaries report a correlation between activation and expression while the title and abstract claim storage (V2), and the same editing machinery moves non-factual linguistic patterns, so the construct never separates from its neighbor (C4) \\
\addlinespace[2pt]
Manipulating these neurons changes how strongly the model expresses the fact
  & Causally Suggestive
  & Specificity, tested and inconclusive (I4); sufficiency reaches a 34.4\% edit success rate on a fact the rest of the model still expresses (I2) \\
\addlinespace[2pt]
Editing these neurons edits that fact and leaves unrelated knowledge alone
  & Disconfirmed
  & Nothing --- tested and failed: Table 6 gives an inter-relation perplexity rise of 7.2 for the identified neurons against 4.3 for random ones, and \S5.1 reads the same table as little negative influence on other knowledge (I4) \\
\addlinespace[2pt]
The account holds beyond BERT-base-cased
  & Insufficient
  & One model, and the generalization is asserted with no experiment behind it (E4, V5); what has since transferred is the attribution procedure rather than the storage reading \\
\bottomrule
\end{tabular}
\end{table}

\begin{table}[H]
\centering
\footnotesize
\caption{\textbf{Knowledge neurons.} Verdict: \textbf{Disconfirmed}. A dash means the tier is reached; a criterion at
Inconclusive, Untested or Disconfirmed blocks promotion.}
\label{tab:knowledge_neurons-verdict}
\renewcommand{\arraystretch}{1.18}
\begin{tabular}{@{}l>{\raggedright\arraybackslash}p{3.6cm}>{\raggedright\arraybackslash}p{4.4cm}@{}}
\toprule
\textbf{Tier} & \textbf{Requires} & \textbf{Missing} \\
\midrule
Proposed & C1--C2; one admissible measurement & --- \\
\addlinespace[2pt]
Causally Suggestive & Necessity (I1); baselines (M2) & --- \\
\addlinespace[2pt]
Mechanistically Supported & Sufficiency (I2); reach (E1); specificity (I4) & \textbf{I4} inconclusive: erasure spares other relations, while updating damages them more than editing random neurons does \\
\addlinespace[2pt]
Triangulated & Convergence (C3); replication (E2/E4); dissociation (I6) & \textbf{I6} unattempted: no second mechanism is shown intact under knowledge-neuron ablation; \textbf{C3}: one attribution family against one contrast baseline \\
\addlinespace[2pt]
Validated & Calibration (M1--M7); interpretive audit (V1--V5); reach (E2, E4) & \textbf{V2} Disconfirmed; \textbf{V3} inconclusive; \textbf{M1}, \textbf{M3}, \textbf{M5}, \textbf{M7} unattempted, and the neuron count is set by a target range rather than by the data \\
\bottomrule
\end{tabular}
\end{table}


\end{document}
